% !TEX root = critical_path_evolution_ru.tex
\documentclass[12pt,a4paper]{article}
\usepackage{cmap}
\usepackage[margin=1in]{geometry}
\usepackage{amsmath,amssymb,amsfonts}
\usepackage[utf8]{inputenc}
\usepackage[T2A]{fontenc}
\usepackage[russian]{babel}
\usepackage[hidelinks]{hyperref}
\usepackage{comfortaa}
\renewcommand*\familydefault{\sfdefault}
\usepackage[x11names]{xcolor}
\usepackage{indentfirst}

\makeatletter
\everydisplay{\global\@afterindenttrue}
\appto\endequation{\global\@afterindenttrue}
\makeatother

\makeatletter
\renewcommand{\thefootnote}{¹}
\renewcommand{\@makefnmark}{\mbox{\textcolor{white}{\thefootnote}}}
\renewcommand{\@makefntext}[1]{\parindent 1.8em\noindent\mbox{}\llap{\thefootnote}\textcolor{white}{#1}}
\renewcommand{\footnoterule}{\kern-3pt\color{white}\hrule width 0.4\textwidth height 0.4pt\kern 2.6pt}
\makeatother

\title{\textbf{Критический путь эволюции}}
\author{%
  Олег Понфиленок\\[0.35em]%
  {\normalsize Независимый исследователь; \href{mailto:ponfil@gmail.com}{ponfil@gmail.com}}%
}
\date{}

\begin{document}

\pagecolor{black}
\color{white}

\maketitle

\begin{abstract}
Эволюция жизни рассматривается как прохождение $N \sim 10^{14}$ малых высоковероятных шагов усложнения со средним временем шага $\mu$~--- критическим путём от теплового гейта реликтового излучения до технологической цивилизации за время $T \approx 13{,}8$~Gyr. При независимости времён шагов концентрация их суммы сжимает относительный разброс времён выхода как $1/\sqrt{N}$; в рамках модели отсюда следует дисперсия $\sigma = \sqrt{\mu T}$. Мы показываем, что наблюдаемое отсутствие «громких» цивилизаций (космическая тишина) даёт \emph{верхнюю границу} на межпланетную дисперсию времён выхода $\sigma \approx 800$--$1400$~лет~--- порядка горизонта перехода к детектируемости; из этой $\sigma$ следует $\mu \approx 25$--$75$~мин, что совпадает с генерационным временем микроорганизмов. Та же тишина \emph{ограничивает сверху} населённость Галактики $n \sim 10^4$--$10^5$ цивилизационных планет, согласуясь с астрофизическими оценками числа землеподобных местообитаний. Два свойства критического пути~--- протекание при локальном избытке ресурсов и кинематическая уникальность самого быстрого маршрута~--- делают его параметры приближённо универсальными и обеспечивают синхронный темп развития жизни; космическая тишина указывает, что этот маршрут реализуется через углеродно-органическую химию жизни нашего типа. Логическое следствие той же картины~--- конвергенция формы цивилизации: носители технологических цивилизаций ожидаемо близки по \emph{функциональному} плану к гуманоидам, хотя детали морфологии могут различаться.
\end{abstract}

\noindent\textbf{Ключевые слова:} парадокс~Ферми; критический~путь~эволюции; модель~простых~шагов (против «трудных~шагов»); синхронизация~времён~появления~разума; MEPP; конвергентная~эволюция; населённость~Галактики; SETI

%%% ===================================================================
\section{Введение}
%%% ===================================================================

Парадокс Ферми~--- отсутствие наблюдаемых техносигнатур при ожидаемой обитаемости галактики~--- обычно обсуждается через редкие «трудные шаги» эволюции (модель Картера~\cite{Carter1983}) или через антропные аргументы о ранней позиции наблюдателя (модель grabby aliens~\cite{Hanson2021}). В данной работе предлагается иной, строго копернианский подход: эволюция жизни от рождения Вселенной до технологической цивилизации идёт по \emph{критическому пути}~--- лимитирующему, самому быстрому маршруту усложнения (по аналогии с методом критического пути в управлении проектами, critical-path method), а не по набору редких маловероятных переходов. Это \textbf{противоположно} понятию «critical steps» Картера: critical steps~--- мало и они трудны; critical path~--- много простых шагов, и важна лишь скорость лимитирующего маршрута.

Выбор в пользу простых шагов мотивирован буквальным прочтением Дарвиновского \emph{natura non facit saltum}~\cite{Darwin1859}: природа не делает скачков, усложнение~--- цепочка малых изменений, а не редких «трудных» переходов. Дополнительный стимул~--- асимметрия эмпирических данных: гипотеза «трудных шагов» сама нигде не доказана~--- неразложимость не была строго показана ни для одного предполагаемого кандидата,~--- тогда как при детальном изучении конкретные кандидаты (эукариогенез, происхождение жизни) раз за разом распадаются на цепочки кооперативных, высоковероятных подшагов~\cite{Mills2025}. Единственность происхождения сама по себе трудность не доказывает: возможны эффект инкумбента (первая успешная линия исключает повторные попытки), закрытие ниши средой и недонаблюдаемость вымерших альтернативных линий.

Термодинамически это можно понимать как отбор наиболее диссипативного маршрута: в открытой системе при избытке локальных ресурсов принцип максимального производства энтропии (MEPP~\cite{Martyushev,MartyushevSeleznev2006,Kleidon2010,Vallino2010}) выделяет путь, быстрее всего превращающий доступную свободную энергию в устойчивое усложнение. Поэтому критический путь~--- не только кинематически самый быстрый, но и физически наиболее диссипативный маршрут. При этом усложнение в избытке ресурсов идёт параллельно по многим допустимым маршрутам, ограниченным физикой и химией; MEPP и кинематический минимум полного времени выделяют среди них единственный самый быстрый~--- критический путь. Синхронизация и конвергенция формы относятся только к его ведущей ветви, а не ко всему филогенетическому дереву (\S~\ref{sec:time_form}).

Эволюцию от теплового гейта реликтового излучения до технологической цивилизации мы рассматриваем как \emph{единую} цепочку $N$ очень многих простых шагов. При сложении многих независимых шагов среднее время растёт как $N$, а разброс~--- лишь как $\sqrt{N}$, поэтому \emph{относительный} разброс времён выхода сжимается с числом шагов: разные планеты приходят к цивилизации почти одновременно~--- в этом и состоит \textbf{синхронизация}. Она держится на двух свойствах критического пути: он \emph{единствен} и идёт с \emph{универсальным} тактом $\mu$ на плато скорости (при избытке ресурсов темп задаётся химией, а не средой). Количественный вывод этих свойств дан в \S~\ref{sec:steps}, порядковая оценка $\sigma$ из тишины~--- в \S~\ref{sec:inversion}; главная цель~--- показать качественный механизм синхронизации, а не точную калибровку временных масштабов.

По сути это иной взгляд на эволюцию~--- не как на череду случайностей, а как на закономерность: отдельный шаг случаен, но концентрация суммы огромного их числа делает итоговый путь практически детерминированным, и жизнь предстаёт не лотереей, а \emph{проектом} с единственным самым быстрым маршрутом усложнения. Поскольку этот маршрут кинематически уникален, модель предполагает как следствие конвергенцию не только времён выхода, но и \emph{формы цивилизации}: технологические цивилизации на других планетах \emph{ожидаемо} близки по \emph{функциональному} плану к гуманоидам (\S~\ref{sec:time_form}).

Ближайший союзник по качественному тезису~--- работа Mills et al.~\cite{Mills2025} («нет трудных шагов»); вопрос о дисперсии времён появления разума ранее ставил Hair~\cite{Hair2011}, но с противоположным выводом: он подчёркивал \emph{большой} разброс и колоссальное временное преимущество древнейших цивилизаций, тогда как из суммы независимых шагов мы получаем узкий синхронный пик. Настоящая работа разделяет вывод «трудных шагов нет», но добавляет количественную рамку суммы последовательных шагов и оценку $\sigma$ из тишины (\S~\ref{sec:inversion}). От модели grabby aliens~\cite{Hanson2021} мы отличаемся отказом от степенного закона hard steps и от антропной посылки: синхронность выводится из концентрации суммы независимых шагов (аддитивность дисперсии) и локальной конвергентности космической фазы. Иную трактовку тишины дают Sandberg, Drexler и Ord~\cite{Sandberg2018}: парадокс «растворяется» статистически~--- при честном учёте многопорядковой неопределённости параметров уравнения Дрейка вероятность одиночества велика. Но главный источник этой неопределённости~--- допущение о возможном крайне маловероятном (трудношаговом) абиогенезе; наш отказ от трудных шагов обрезает этот хвост, неопределённость смещается к «жизнь не редка», и необходимость объяснять тишину возвращается~--- мы отвечаем на неё синхронностью. Наблюдательно модели различимы формой распределения времён выхода~--- узкий синхронный пик против широкого хвоста «трудных шагов».

%%% ===================================================================
\section{Критический путь из модели простых шагов}
\label{sec:steps}
%%% ===================================================================

\textit{Данный раздел даёт количественный вывод: формализует дисперсию времён выхода как концентрацию суммы большого числа независимых шагов и обосновывает универсальность параметров критического пути.}

\subsection{Допущения и концентрация суммы шагов}
\label{sec:concentration}

Постулат: эволюционное усложнение идёт $N$ малыми высоковероятными шагами (по Дарвину~\cite{Darwin1859}, \emph{natura non facit saltum}; не редкими «трудными», \S~\ref{sec:planetary_steps}). Время выхода $T_{\mathrm{exit}} = \sum_{i=1}^{N} t_i$, шаги независимы. Отдельные шаги \emph{не} обязаны быть одинаковыми по длительности: $i$-й имеет своё среднее $\mu_i$ и дисперсию $\sigma_i^2$; $\mu \equiv \langle t_i\rangle = T/N$~--- \emph{среднее} время шага по цепочке, $\sigma_{\mathrm{s}}^2 \equiv \langle\sigma_i^2\rangle$~--- средняя пошаговая дисперсия (\S~\ref{sec:heterogeneity}).

Ядро механизма~--- \textbf{аддитивность среднего и дисперсии} (тождество Бьенеме): $\mathbb{E}[T_{\mathrm{exit}}] = \mu N$, $\mathrm{Var}(T_{\mathrm{exit}}) = N\sigma_{\mathrm{s}}^2$, откуда $\sigma = \sigma_{\mathrm{s}}\sqrt{N}$ и относительный разброс $\sigma/T = c_v/\sqrt{N}$ при $c_v = \sigma_{\mathrm{s}}/\mu$. Для самой концентрации достаточно независимости шагов и конечности $c_v$; нормальность одношаговых распределений не требуется (нормальность суммы~--- лишь в \S~\ref{sec:silence_sigma}). Частный случай $c_v \approx 1$ даёт $\sigma = \sqrt{\mu T}$; на нём держится числовая инверсия в $\mu$ (\S~\ref{sec:inversion}), но не сам вывод о концентрации. Влияние неоднородности $\mu_i$ на $\sigma$ сводится к множителю порядка единицы, если ни один шаг не доминирует (\S~\ref{sec:heterogeneity}).

Концентрация опирается на два слабых требования: (1)~велико число шагов $N$; (2)~ни один шаг не доминирует в $\sum_i\sigma_i^2$~--- нет неразложимого «трудного» шага с временем порядка целевой $\sigma$ (\S~\ref{sec:planetary_steps}). В опорной форме $\sigma = \sqrt{\mu T}$ при $T \approx 13{,}8$~Gyr~--- полном времени от теплового гейта CMB (\S~\ref{sec:common_start}); из тишины калибруются численные $\sigma$ (\S~\ref{sec:silence_sigma}) и $\mu$ (\S~\ref{sec:mu_from_sigma}).

\subsection{Чувствительность дисперсии к неоднородности шагов}
\label{sec:heterogeneity}

При независимости шагов дисперсия времени выхода равна сумме пошаговых дисперсий:
\begin{equation}
\sigma^2 = \mathrm{Var}(T_{\mathrm{exit}}) = \sum_{i=1}^{N}\sigma_i^2,
\label{eq:varsum}
\end{equation}

Для анализа неоднородности введём два параметра. У шага $i$ среднее $\mu_i$ и внутришаговый разброс $\sigma_i = c_v\mu_i$ с общим $c_v = \sigma_{\mathrm{s}}/\mu$; средние $\mu_i$ вдоль цепочки имеют среднее $\mu \equiv T/N$ и коэффициент вариации $CV_\mu \equiv \mathrm{std}(\mu_i)/\mu$ (относительный разброс длительностей шагов). Из~\eqref{eq:varsum} следует:
\begin{equation}
\boxed{\;\sigma = c_v\,\sqrt{1+CV_\mu^{2}}\;\sqrt{\mu\,T}\;}
\label{eq:sigma_general}
\end{equation}
В~\eqref{eq:sigma_general} разделены два вклада: стохастика внутри шага ($c_v$, \S~\ref{sec:inversion}) умножает $\sigma$ линейно; разброс $\mu_i$ между шагами ($CV_\mu$) входит как множитель $\sqrt{1+CV_\mu^2}$. Относительно опорной $\sigma_0=c_v\sqrt{\mu T}$ (случай равных $\mu_i$, $CV_\mu=0$): при $CV_\mu=1$, когда типичное отклонение длительности шага сопоставимо с его средним, $\sigma/\sigma_0=\sqrt{2}\approx1{,}4$. Зависимость сублинейна по $CV_\mu$, поэтому умеренный разброс длительностей быстрых шагов лишь слабо сдвигает $\sigma$.

Существеннее вклад самого медленного шага. Если $\mu_{\max}\gg\mu$, то
\begin{equation}
\sigma \approx \sqrt{\big(\sqrt{\mu T}\,\big)^2 + (c_v\mu_{\max})^2}.
\end{equation}
Сумма $N\sim10^{14}$ быстрых шагов даёт пол $\sim10^3$~лет (\S~\ref{sec:mu_from_sigma}); медленный шаг добавляется в квадратурах. Доминирование наступает при $c_v\mu_{\max}\gtrsim\sigma\approx10^3$~лет, то есть $\mu_{\max}\gtrsim\sqrt{N}\,\mu\approx10^{7}\mu$: отдельный шаг может быть в $\sim10^{7}$ раз медленнее среднего ($40$~мин $\to$ до $\sim10^3$~лет), не доминируя в сумме; при $\mu_{\max}\gtrsim10^3$~лет он определяет $\sigma$. Эукариогенез как единый неразложимый шаг ($\mu_{\max}\sim2$~Gyr) дал бы $\sigma\sim2\times10^9$~лет~--- на шесть порядков выше наблюдаемой межпланетной дисперсии (\S~\ref{sec:eukaryo}).

Межпланетная $\sigma\approx10^3$~лет задаётся прежде всего числом шагов $N$ и временем самого медленного шага $\mu_{\max}$; разброс длительностей быстрых шагов ($CV_\mu$) вносит лишь слабый множитель. Условие концентрации~--- отсутствие неразложимого «трудного» шага с $\mu_{\max}\gtrsim10^3$~лет; его проверка для планетной фазы~--- в \S~\ref{sec:planetary_steps}.

\subsection{Разложение трудных шагов на простые}
\label{sec:planetary_steps}

На планетарном участке остаётся следующий риск. Длительные планетные эпохи (например, $\sim 2$~млрд лет до эукариот) сами по себе \textbf{не аномальны}: это просто \emph{длинные цепочки} обычных тактов (большое число шагов $N$), дающие тот же закон $\sigma = \sqrt{\mu T}$. Аномальный вклад в межпланетный разброс возник бы лишь от \textbf{неразложимого} медленного планетного шага~--- единичного перехода с собственным временем $\gtrsim10^3$~лет (порог доминирования из \S~\ref{sec:heterogeneity}), который нельзя представить цепочкой высоковероятных подшагов и который не является общим для планет.

Это, в иной форме,~--- гипотеза «трудных шагов» Картера~\cite{Carter1983}: редкие, по-настоящему неразложимые переходы с большой дисперсией времени ожидания. Но за более чем 40 лет обсуждения ни один конкретный шаг не доказан неразложимым; сама гипотеза опирается на допущения, которые недавний пересмотр называет сомнительно обоснованными как антропно, так и эволюционно~--- эффект «убранной лестницы» (успешные линии вытесняют конкурентов, маскируя их независимое происхождение), вытеснение нишы изменением среды и принципиальная ненаблюдаемость вымерших независимых линий~\cite{Mills2025}. Текущая тенденция исследований прямо противоположна: каждый из кандидатов на «трудный шаг» при ближайшем рассмотрении распадается на цепочку простых~--- происхождение эукариот через архей группы Asgard и их культивируемого представителя \emph{Prometheoarchaeum syntrophicum}~\cite{Spang2015,Zaremba2017,Imachi2020}, синтрофные и водородные гипотезы симбиогенеза~\cite{Martin1998,LopezGarcia2020}, актиновые гомологи у Lokiarchaeota~\cite{Spang2015}, наблюдаемый актиновый цитоскелет у Asgard-археи~\cite{Rodrigues2023} и новые филогеномные реконструкции 2025--2026~гг.~\cite{Tobiasson2026,Speijer2025,BravoArevalo2025}.

Таким образом, вся нагрузка ложится на разложимость планетных шагов: пока не предъявлен ни один доказанно неразложимый шаг, постулат модели имеет эмпирическую поддержку, но не доказательство. Тест эукариогенеза (\S~\ref{sec:eukaryo}) остаётся решающим именно потому, что это последний серьёзный кандидат на «трудный шаг», для которого общепринятой механистической модели пока нет.

\subsection{Каскадная модель: направленность и рост диссипации}
\label{sec:cascade}

Усложнение~--- каскад переходов по энергетическим ступенькам со свободной энергией Гельмгольца $\Delta F = \Delta U - T\Delta S$. Каждая ступенька~--- один шаг~--- метастабильное состояние с памятью (время удержания $\tau \sim \exp(\Delta U/kT)$); производство энтропии эффективно понижает барьеры, делая каскадный подъём проходимым.

\textbf{Направленность} возникает из асимметрии: с ростом сложности система растёт $\Rightarrow$ больше степеней свободы $\Rightarrow$ плотность микросостояний увеличивается $\Rightarrow$ переход «вправо» энтропийно выгоднее (фазовый объём более сложных состояний больше). В большей степени асимметрия проявляется в \textbf{высоте} шага, чем в его \emph{темпе}: каждый шаг вправо сопровождается бóльшим приростом энтропии $\Delta S$ и бóльшей диссипацией свободной энергии. Вместе с репликацией (храповик, \S~\ref{sec:ratchet}) это даёт статистический дрейф к усложнению вдоль стрелы времени (телеономия; ср.~Lotka~\cite{Lotka1922}, MEPP~\cite{Martyushev}, классическое термодинамическое обоснование~\cite{SchneiderKay1994}, диссипативная адаптация в самосборке~\cite{England2015}).

Производство энтропии ускоряется «вправо» по двум причинам: (i) \textbf{экспансия}~--- репликация наращивает число систем на переднем фронте экспоненциально; (ii) более сложные метастабильные структуры требуют большего удельного притока свободной энергии. Количественно эта тенденция фиксируется независимо измеряемой \emph{плотностью потока свободной энергии} $\Phi_m$ (эрг$\cdot$с$^{-1}\cdot$г$^{-1}$), которая растёт на порядки вдоль лестницы: звезда $\sim1$, планета $\sim10^2$, растение $\sim10^3$, животное $\sim10^4$, мозг $\sim10^5$, технологическое общество $\sim10^6$~\cite{Chaisson2011}. Само же \emph{среднее} время элементарного такта $\mu$ (один цикл репликации, один акт наследуемой памяти) снизу ограничено биохимией копирования наследуемой памяти~--- поэтому $\mu \approx$ время генерации (\S~\ref{sec:resource_threshold}). Систематическое различие \emph{средних} тактов между крупными фазами~--- абиотической, одноклеточной и многоклеточной (\S~\ref{sec:sex})~--- не отменяет нижней границы $\mu$ на плато микробной скорости, но смещает путевую оценку из однофазной инверсии в пределах фактора в несколько раз.

\subsection{Биологический храповик}
\label{sec:ratchet}

В одиночной механической системе видны и прямые, и обратные переходы. В биологии обратное движение (вымирание, редуктивная эволюция) идёт постоянно, но на уровне отдельных линий, и эти линии исчезают, не оставляя наблюдателя.

Ключевой усилитель~--- \textbf{репликация с памятью}: удачная ступенька копируется, и число систем на ней растёт экспоненциально. Чтобы откатилась вся популяция, независимый откат должен случиться у огромного числа копий разом~--- вероятность пренебрежимо мала. Репликация превращает слабую термодинамическую асимметрию в практически необратимый \textbf{храповик}~--- передний фронт сложности назад не откатывается (ср.~биологический закон необратимости Долло~\cite{Dollo1893,Gould1970}). Термодинамически этот храповик можно проиллюстрировать концепцией Джереми Ингланда: деление бактерии на две части~--- процесс относительно простой, но спонтанное слияние двух образовавшихся половин обратно в одну функциональную бактерию статистически и термодинамически невероятно из-за диссипации тепла и огромной разницы в энтропии состояний~\cite{England2013}.

\subsection{Плато скорости и ресурсный порог}
\label{sec:resource_threshold}

Два условия универсального такта $\mu$ на микробных шагах~--- \emph{плато скорости} адаптации и ничтожный \emph{ресурсный порог}~--- обосновываются ниже популяционно-генетически и энергетически.

В бесполой популяции одновременно конкурируют многие полезные мутации: пока одна линия фиксируется, сопоставимые по приспособленности линии отбрасываются~--- это \emph{клональная интерференция}. Скорость адаптации растёт с численностью $N_c$ и потоком полезных мутаций $U_b$ лишь \textbf{логарифмически} и при больших $N_c$ \textbf{насыщается} (модель Герриша--Ленски~\cite{GerrishLenski1998}, travelling-wave теория~\cite{DesaiFisher2007,ParkKrug2007}; экспериментально~--- «предел скорости», не зависящий от потока мутаций~\cite{deVisser1999,Lang2013}). Насыщение наступает при $N_c\,U_b \gtrsim 1$, т.е.\ при $N_c \gtrsim 10^{6}$ клеток ($U_b \sim 10^{-6}$); дальнейший рост $N_c$ скорость почти не меняет~--- приспособленность не успевает передаваться через последовательные выметания. Снизу её ограничивает \emph{биохимический пол}~--- минимальное время деления ($\mu_{\min}\approx$ время генерации, $\sim 20$~мин у \emph{E.~coli}; такт прижат именно к этому полу, поскольку критический путь идёт по самой быстрой линии~--- первому успеху среди множества параллельных попыток).

Для оценки энергобюджета критического пути берём \textbf{эволюционный фронт максимального темпа}~--- популяцию, уже находящуюся на этом плато. Эталон~--- одна колба долгосрочного эксперимента Ленски (LTEE~\cite{Lenski2017}): 10~мл, $\sim 5 \times 10^{8}$ клеток при суточном пике численности ($\sim 5 \times 10^{7}$~клеток/мл). Это на $\sim 2$--$3$ порядка \emph{выше} порога $N_c \sim 10^{6}$; эффективный размер $N_e \approx 3 \times 10^{7}$~--- тоже глубоко на плато. Именно поэтому LTEE наблюдает максимальный кинетический темп~--- все 12 линий идут почти одной кривой приспособленности~\cite{Wiser2013,Lenski1994},~--- и служит правильным масштабом «переднего фронта». Его энергетика (по $\sim 0{,}57$~pW на клетку в экспоненциальной фазе~\cite{Deng2021ATP}):
\begin{itemize}
\item мощность $\approx 3 \times 10^{-4}$~Вт; $\dot{S}_{\mathrm{front}} \approx 10^{-6}$~Вт/К.
\end{itemize}
Планета (Земля): $\dot{S}_{\mathrm{planet}} \approx 4{,}6 \times 10^{14}$~Вт/К; $\dot{S}_{\mathrm{bio}} \approx 4 \times 10^{11}$~Вт/К; прокариот $\sim 5 \times 10^{30}$.

Отношение фронт/планета: $\dot{S}_{\mathrm{front}}/\dot{S}_{\mathrm{planet}} \approx 2 \times 10^{-21}$; клетки $\approx 10^{-22}$; площадь~--- фронту достаточно $\sim 2$~см$^2$ ($\sim 3 \times 10^{-19}$ поверхности планеты). Это снимает \emph{энергетическое и ресурсное} лимитирование пути глобальными условиями планеты.

\subsection{Локальные оазисы эволюции на планетах}
\label{sec:applicability_bounds}

Концепция критического пути вступает в содержательную дискуссию с гипотезой «глобального средового гейтирования» (environmental gating), отстаиваемой, в частности, в работе Mills et al.~\cite{Mills2025}. Согласно их взглядам, крупные эволюционные переходы (эукариогенез, многоклеточность) задерживались не из-за внутренней сложности, а в ожидании созревания планеты (достижения глобального порога кислорода, климатической стабилизации и т.д.). Такой механизм неизбежно приводил бы к значительному разбросу времён выхода цивилизаций, поскольку планетные таймеры (скорость накопления $O_2$) у всех планет уникальны.

В рамках нашей модели геология и планетарные таймеры не являются сдерживающими шлюзами (гейтами), благодаря двум эмпирически подтвержденным факторам:
\begin{enumerate}
   \item \textbf{Кислородные оазисы.} Биохимическая эволюция не ждет глобального изменения планеты, а протекает в локальных очагах с избытком ресурсов. Геологические данные показывают, что мелководные «кислородные оазисы», создаваемые локальными сообществами цианобактерий, существовали сотни миллионов лет до Великого окислительного события (GOE) в условиях полностью аноксичной планеты~\cite{Riding2014,Ostrander2019}. Эволюционные шаги, требующие кислорода, шли внутри этих локальных оазисов на максимальной кинетической скорости, а последующее глобальное насыщение планеты привело лишь к пространственному расселению уже сформированных форм жизни.
   \item \textbf{Тест порядка появления (Mills~2014 против Mills~2025).}  Гипотеза гейтирования утверждает, что появление животных жестко лимитировалось ростом кислорода. Однако независимые измерения потребности ранних многоклеточных (на примере современных губок) показывают, что они способны выживать при концентрациях кислорода всего $0{,}5$--$4{,}0\%$ от современного уровня~\cite{Mills2014}. Такие концентрации были достигнуты на Земле задолго до реального появления животных. Тот факт, что животные возникли со значительным временным лагом после прохождения этого порога, свидетельствует, что кислород не был лимитирующим стоп-фактором. Эволюция шла со своим собственным кинетическим темпом биологических переходов.
\end{enumerate}

\label{sec:plateau}
На плато скорости такт шага задаётся химией ($\mu_{\min,i}$ для шага $i$), а не средой, и одинаков на всех планетах (\S~\ref{sec:planetary_spread}); порог ресурсов ничтожен (\S~\ref{sec:resource_threshold}), эволюция идёт в локальных оазисах, а не в ожидании глобального созревания планеты. Планетам не нужно быть идентичными: достаточно, чтобы в \emph{каждый момент} была хотя бы одна ниша на плато и возможность быстрой экспансии между оазисами; разные стадии могут требовать разных мест, поэтому критический путь опирается на колонизацию и миграцию между оазисами.

По мере движения «вправо» по каскаду усложнения (\S~\ref{sec:cascade}) объём потребляемых ресурсов и энергопотребление монотонно растут (растёт высота ступеньки и производство энтропии), а кажущийся темп эволюции ускоряется за счёт экспансии. Поэтому тезис «локального оазиса достаточно» мы относим прежде всего к ранним шагам с малым энергобюджетом фронта относительно планеты (\S~\ref{sec:resource_threshold}); количественная привязка к шкале Кардашёва здесь не делалась. Уже на нашем уровне ($K \approx 0{,}73$) цивилизация требует существенной доли планетарного энергетического пула; дальше потребление будет только нарастать (\S~\ref{sec:pred_kardashev}).

Для \emph{поздних} шагов, требующих глобальных условий, планетный фактор может оставаться значимым: например, возникновение крупных активных хищников с высоким метаболизмом, которые физически не могут существовать в границах локального оазиса~\cite{Knoll2014}. В таких редких звеньях критический путь переходит на планетный масштаб~--- это естественная граница применимости модели: такт $\mu = \mu_{\min}$ на плато скорости и независимость от глобальной среды хорошо обоснованы для подавляющего большинства шагов и времени пути, но не для всех поздних переходов без оговорок.

\subsection{Подавление планетарной общей моды}
\label{sec:planetary_spread}

Опорная форма $\sigma = \sqrt{\mu T}$ (\S~\ref{sec:concentration}) законна лишь
тогда, когда межпланетный разброс времён выхода исчерпывается \emph{независимой}
пошаговой стохастикой. По закону полной дисперсии, если средний такт плато
$\mu_{\min}$ сам варьирует между планетами с относительным разбросом
$\delta \equiv \mathrm{std}_{\mathrm{пл}}(\mu_{\min})/\mu_{\min}$, то
\begin{equation}
\sigma^2_{\mathrm{межпл}}
  = \underbrace{N\sigma_{\mathrm{s}}^2}_{\text{независимая}}
  + \underbrace{T^2\,\delta^2}_{\text{общая мода}}.
\label{eq:total_variance}
\end{equation}
Первый член совпадает с \S~\ref{sec:concentration} и даёт $\sigma = c_v\sqrt{\mu T}\sim10^3$~лет (\S~\ref{sec:silence_sigma}).
Второй~--- общемодовый: он \emph{не} подавляется как $1/\sqrt{N}$ и доминирует уже
при $\delta \gtrsim c_v/\sqrt{N}\sim10^{-7}$, обращая $\sigma \approx T\,\delta$.
Поэтому синхронность требует, чтобы такт плато был межпланетно инвариантен с этой
точностью. Ниже мы показываем, что инвариантность \emph{обеспечивается отбором на
оптимум}, а не постулируется: планетарный разброс входит не как непрерывное
замедление такта, а как бинарный отбор в когорту.

Планетарная «общая мода» разрушала бы концентрацию суммы шагов: она возникла бы,
если межпланетный разброс задавался \emph{условиями планеты} (температурой,
составом, климатом), общими для всех её шагов и потому коррелирующими их между
собой. Критический путь по построению задаётся не ими, а \emph{универсальными
свойствами химических элементов}: он всегда идёт в локальных оазисах при избытке
ресурсов, на плато $\mu = \mu_{\min}$ (\S~\ref{sec:applicability_bounds}). Общепланетного множителя замедления нет; в
формуле~\eqref{eq:total_variance} должна остаться лишь независимая пошаговая
стохастика~--- ровно та, для которой работает $\sigma/T = c_v/\sqrt{N}$.

Три механизма удерживают планетный вклад в $\delta$ ниже порога~\eqref{eq:total_variance}. \textbf{Температура}~--- непрерывное поле: на
планете с температурным градиентом всегда существует место (широта, глубина,
окрестность источника), где $T = T_{\mathrm{opt}}$ \emph{точно}; оазис садится на
максимум скорости независимо от среднего климата, и температурный вклад в $\delta$
\emph{исчезает}, а не подавляется квадратично. \textbf{Внутриклеточный состав}
(pH, ионная сила, концентрации субстратов и кофакторов) нормализуется активным
транспортом и бионакоплением: клетка концентрирует нужные виды до оптимума,
\emph{развязывая} удельный такт от валовой распространённости в резервуаре.
Универсальный нуклеосинтез предшествующих поколений звёзд гарантирует присутствие
органогенных и кофакторных элементов повсюду в диске
(\S~\ref{sec:local_conv})~--- так что концентрировать всегда есть что; при этом
типичный $\sim 2$--$3$-кратный разброс \emph{соотношений} этих элементов в диске
не переносится на внутриклеточный оптимум. \textbf{Катализатор}~--- кинематическая уникальность
(\S~\ref{sec:uniqueness}) фиксирует на каждом лимитирующем шаге не только число
шагов, но и саму \emph{быстрейшую химию}. На плато скорости лимитирующий шаг прижат к
\textbf{физическому потолку}~--- диффузионному или заданному энергетикой
переходного состояния (эмпирически $k_{\mathrm{cat}}/K_m\sim10^{8}$--$10^{9}$~М$^{-1}$с$^{-1}$,
диффузионный предел~\cite{Bar-Even2011}); он определяется константами и химией
реакции, а не планетой, так что $\mu_{\min}$ межпланетно инвариантен в первом
порядке и $\delta$ остаётся ниже порога~\eqref{eq:total_variance}.

Граница применимости~--- та же, что в \S~\ref{sec:resource_threshold}: если на
планете \emph{нет} ни одной ниши, достигающей оптимума в данный момент~---
требуемого элемента нет вовсе (а не просто меньше), нет
жидкофазной среды при $T_{\mathrm{opt}}$, исчерпан ресурсный пул,~--- планета не идёт по критическому
пути на максимальной скорости и \textbf{выпадает из синхронизированной когорты}
(\S~\ref{sec:bimodal}).

\subsection{Единственность критического пути}
\label{sec:uniqueness}

Усложнение в открытой системе идёт параллельно по многим маршрутам~--- всем, до которых можно дотянуться при данных ресурсах и условиях. В принципе допустимые пути задаёт физика и химия Вселенной; на практике их перечень сужают локальная доступность веществ, энергии и условий. Там, где ресурсы избыточны, а спектр условий достаточен (\S~\ref{sec:resource_threshold}, \S~\ref{sec:plateau}), ресурсное ограничение снимается и доступен весь теоретически возможный набор маршрутов. На этом множестве принцип максимального производства энтропии (MEPP~\cite{Martyushev,MartyushevSeleznev2006,Kleidon2010,Vallino2010}) отбирает наиболее диссипативные траектории; поскольку рост сложности и есть рост производства энтропии (\S~\ref{sec:cascade}), самые диссипативные совпадают с самыми быстрыми по усложнению. Среди них существует единственный минимум полного времени~--- точные совпадения скоростей имеют нулевую меру; это и есть \emph{критический путь}.

Здесь важно оговориться, что критический путь по определению~--- не путь, привязанный к конкретной химии, а \textbf{максимум по скорости среди всех возможных маршрутов} усложнения, допускаемых физикой и химией данной Вселенной. До выхода на громкую фазу критический путь может быть только один: если бы существовала химическая основа, реализующая более быстрый маршрут, мы уже наблюдали бы его результат~--- громкую цивилизацию, выросшую на этом более быстром пути. Поскольку громких цивилизаций не наблюдается (космическая тишина), углеродно-органическая химия, на которой развивается жизнь нашего типа, совпадает с критическим путём данной Вселенной.

Теоретически возможны альтернативные формы жизни на иной химической основе, но они не образуют \emph{собственные} критические пути~--- по построению это просто более медленные маршруты, не достигающие статуса критического. Они также могут в конечном итоге приводить к возникновению разума и цивилизаций, но за существенно более длительные масштабы времени. Сам критический путь задаётся не выбором химии, а фундаментальными константами и параметрами Вселенной (скорость света, гравитационная и тонкая структурная постоянные, соотношения масс частиц и т.~д.): более быстрый критический путь в принципе мог бы существовать только в иной Вселенной с другими физическими константами.

\textbf{Кинематическая уникальность фиксирует число шагов ($N = \mathrm{const}$).} То, что выглядит единичным «ключевым переходом» (эукариогенез, многоклеточность и т. п.), на критическом пути есть длинная цепочка таких шагов. Поэтому шагов очень много~--- $N \sim 10^{14}$ (\S~\ref{sec:inversion}),~--- но их число фиксировано. Критический путь определяется не минимумом \emph{числа} шагов, а \textbf{минимумом полного времени}: это самый короткий по времени путь к тому же уровню развития среди всех возможных альтернативных маршрутов~--- от пребиотической химии до заданного порога сложности (мы берём выход на громкую фазу~--- Тип~II по Кардашёву; \S~\ref{sec:loud_time}). В рамках фундаментальных законов физики и химии такой самый быстрый маршрут кинематически уникален и может быть только один. Любые альтернативные маршруты существенно медленнее: на них шагов больше, сами шаги могут отличаться, а суммарное время заметно дольше. Поскольку самый быстрый путь физически уникален~--- задан фундаментальными константами Вселенной, а не выбором химии,~--- число шагов $N$ на нём жёстко зафиксировано и не имеет межпланетного разброса ($\Delta N = 0$).

\textbf{Критический путь~--- одна ветвь, а не всё дерево.} Критический путь~--- это \emph{единственная} ведущая линия, несущая в каждый момент передний фронт сложности и диссипации; ретроспективно её след на Земле~--- это линия предков человека (прокариот $\to$ эукариот $\to$ многоклеточное $\to$ \ldots $\to$ примат $\to$ технологическая цивилизация), а не вся флора и фауна. Важно, что эта линия определяется \emph{не} телеологически (конкретный вид не «задан» заранее), а как маршрут, удерживающий в каждый момент рекорд по скорости производства энтропии; ретроспективно он и совпадает с предками самого диссипативного состояния~--- цивилизации. Подавляющее большинство ветвей филогенетического дерева~--- прочие виды, целые царства~--- лежат \emph{вне} критического пути: это контингентные ответвления (ср. Cit$^+$ в LTEE, \S~\ref{sec:ltee}), которые на разных планетах могут различаться как угодно. Универсальность и синхронность модель утверждает \emph{только} для ведущей ветви: синхронизируется момент выхода цивилизации, а не состав экосистемы. Соответственно и конвергенция формы цивилизации (\S~\ref{sec:time_form}) относится лишь к этой линии~--- сопутствующая биота на другой планете может быть совершенно иной.

\subsection{Критический путь при половом размножении}
\label{sec:sex}

В многоклеточной фазе элементарный \emph{генетический} такт снизу ограничен временем поколения~--- месяцы или годы, а не минуты микробной фазы. Если учитывать только последовательные наследуемые мутации с $\mu_{\mathrm{gen}}\sim1$~год на фоне фазы длительностью $10^8$--$10^9$~лет, то по \S~\ref{sec:heterogeneity} вклад этой фазы в $\sigma$ составляет $\sim\sqrt{T_{\mathrm{phase}}\mu_{\mathrm{gen}}}\sim10^4$~лет~--- на порядок больше наблюдаемой межпланетной дисперсии $\sim10^3$~лет.

Модель относит этот масштаб не к календарному поколению, а к \emph{эффективному} такту $\mu_{\mathrm{eff}}$: полезные изменения возникают в разных особях одновременно, а рекомбинация собирает их в одной линии~\cite{Muller1932,HillRobertson1966,Colegrave2002}. За одно поколение геном может получить сразу несколько накопленных аллелей~--- календарно это один перекрывающийся интервал, а не сумма $n_{\mathrm{gen}}$ последовательных ожиданий; при $n_{\mathrm{gen}}$ параллельно собранных инкрементов $\mu_{\mathrm{eff}}\sim\mu_{\mathrm{gen}}/n_{\mathrm{gen}}$.

Два режима задают порядок $n_{\mathrm{gen}}$. \textbf{Классические свипы} ограничивают \emph{число фиксаций}: рекомбинация не позволяет накапливать больше $\sim$одной замены на сантиморган за $\sim200$ поколений~\cite{WeissmanBarton2012}~--- десятки фиксаций на геном за поколение, что для $n_{\mathrm{gen}}\sim10^2$--$10^3$ недостаточно. \textbf{Полигенная адаптация} действует иначе: сложные признаки сдвигаются согласованным изменением частот на \emph{сотнях--тысячах} локусов из стоячей вариации~\cite{Pritchard2010,Boyle2017,Barghi2020}; популяционное среднее меняется долговечно, хотя отдельные аллели не фиксируются. Этот режим не упирается в потолок свипов и даёт $n_{\mathrm{gen}}\sim10^2$--$10^3$ параллельно изменяемых локусов за поколение~--- столько, чтобы $\mu_{\mathrm{eff}}$ оказался на два порядка короче $\mu_{\mathrm{gen}}$ и вклад фазы в $\sigma$ имел порядок $\sim10^3$~лет. Так критический путь в фанерозое идёт мелкими параллельными шагами (\emph{natura non facit saltum}, \S~\ref{sec:concentration}), а не редкими свипами; храповик здесь~--- на уровне \emph{популяционного среднего} фенотипа. Часть субгенерационных инкрементов (обучение, культура, эффект Болдуина~\cite{Baldwin1896}, \S~\ref{sec:cascade}) переводит количественные сдвиги в более дискретные наследуемые ступеньки. Количественная доля таких каналов в фанерозое~--- открытый вопрос (\S~\ref{sec:macro_tempo}).

\subsection{Современные шаги~--- технологические, а не генетические}
\label{sec:modern_steps}

На нынешнем этапе критический путь движется не медленной биологической эволюцией человека, а цивилизационно-технологическим прогрессом. При этом интернет, компьютер, авиация, искусственный интеллект~--- это не отдельные шаги, а \emph{вехи}: каждая сама складывается из огромного числа элементарных шагов усложнения (научных работ, инженерных улучшений, актов обмена знанием) и не возникает за один атомарный акт. Атомарный шаг на этом уровне намного мельче такой вехи. Кажущееся ускорение прогресса на поздних стадиях обеспечивается не сокращением такта отдельного шага, а сменой носителя памяти (гены $\to$ культура $\to$ техника).

\subsection{Конвергентность цивилизаций}
\label{sec:time_form}

Кинематическая уникальность (\S~\ref{sec:uniqueness}) относится к \emph{полному времени} $T_{\mathrm{exit}}$: среди всех траекторий к порогу сложности существует единственный глобальный минимум времени. Проекция на пространство \emph{форм цивилизации}~--- морфология биологического носителя, функциональный план, техносфера~--- высокоразмерна, но концентрация суммы шагов сжимает межпланетный разброс по времени. Времена выхода попадают в узкую «трубку» $\sigma \sim 10^3$~лет (\S~\ref{sec:concentration}, \S~\ref{sec:inversion}), а относительный разброс $\sigma/T \sim c_v/\sqrt{N} \sim 10^{-7}$ при $N \sim 10^{14}$. Логично как следствие предположить, что аналогичным образом сжимается и форма: цивилизации на разных планетах, конечно, не тождественны, но их морфологически-функциональный разброс того же порядка малости~--- каждая своя, но в целом все \emph{относительно} близки, \emph{в окрестности} одного гуманоида \emph{функционального типа}.

Из модели (вместе с конвергентной эволюцией, \S~\ref{sec:ltee}) \emph{следует} воспроизводимый \emph{функциональный} план биологического носителя~--- манипуляторные конечности, крупный социальный мозг, орудийность,~--- а не принципиально иной эволюционно-морфологический базис. Остаточные отличия между планетами лежат на масштабе той же «дисперсии»: не произвольный разброс форм, а малые флуктуации вокруг аттрактора критического пути (без утверждения тождества с земной морфологией; \S~\ref{sec:open_civ_form}).

\subsection{Экспериментальное свидетельство воспроизводимости пути}
\label{sec:ltee}

Долгосрочный эксперимент Ленски (LTEE)~\cite{Lenski2017}~--- лучший доступный полигон: 12 параллельных «проигрываний плёнки» из одного предкового клона, более 73\,000 поколений. Это контролируемая постановка спора Конвей-Морриса~\cite{ConwayMorris2003} $\leftrightarrow$ Гулда~\cite{Gould1989}.

Показатель синхронности~--- не идентичность мутаций, а \textbf{повторяемость функционального результата и времени} между независимыми линиями. Все 12 линий идут почти одной кривой роста приспособленности~\cite{Wiser2013,Lenski1994}. На уровне результата эволюция повторяема~--- прямое свидетельство того, что самый быстрый критический путь усложнения фиксирован и воспроизводим.

Контингентные сингулярности существуют, но вне критического пути: аэробный рост на цитрате (Cit$^+$) возник лишь в 1 из 12 линий~\cite{Blount2008,Blount2012}. Это медленный детур к побочному ресурсу; критический путь обходит такие ответвления. Cit$^+$ показывает, что контингентные переходы существуют, но не лежат на критическом пути, пока есть более быстрый маршрут.

\textbf{Оговорка о масштабе.} LTEE меряет микроэволюцию; перенос микро$\to$макро~--- отдельное допущение.

%%% ===================================================================
\section{Эволюция до Земли: космическая фаза критического пути}
\label{sec:cosmic_ensemble}
%%% ===================================================================

\textit{Данный раздел описывает космическую фазу критического пути: синхронный старт от реликтового гейта, три независимые «часы сложности», локально конвергентную пребиотику до порога C*, «заморозку» открытокосмического усложнения выше C* и формирование синхронизированной планетарной когорты.}

\subsection{Общий старт и реликтовый гейт}
\label{sec:common_start}

Критический путь начинается не с появления жизни на планете и не с образования Земли, а с первого момента, когда жидкостная химия стала возможна во Вселенной. Этот момент задаётся \textbf{тепловым гейтом реликтового излучения}: при $100 \lesssim (1+z) \lesssim 137$ температура CMB составляла 273--373~K, делая жидкостную химию возможной на любых каменистых телах независимо от наличия звезды, уже через $\sim 10$--17~млн лет после Большого взрыва~\cite{Loeb2014}. Реликт изотропен до $\Delta T/T \sim 10^{-5}$; при $T_{\mathrm{CMB}} \propto t^{-2/3}$ (эра доминирования материи) разброс момента открытия гейта $\Delta t/t \approx (3/2)(\Delta T/T) \sim 1{,}5 \times 10^{-5}$, откуда при $t \sim 10$--17~млн~лет $\Delta t \sim 10^2$~лет~--- на порядок меньше $\sigma \approx 10^3$~лет, независимо полученной из тишины (\S~\ref{sec:inversion}). Это космологический, а не биологический \emph{общий старт}: «часы» запущены одновременно для всей Вселенной, а в масштабах Галактики ($\sim 30$~кпк, или $\sim 0{,}3$~кпк при $z\sim 120$) шелковское затухание~\cite{Silk1968} подавляет флуктуации CMB до $\ll 10^{-5}$, так что $\sigma_{\mathrm{CMB}} \approx 0$ и межсистемный разброс от гейта отсутствует. Межпланетную дисперсию определяет лишь стохастика суммы шагов (аддитивность дисперсии) по всей цепочке усложнения~--- от этого момента до технологической цивилизации. Какой именно момент считать общим стартом~--- тепловой гейт CMB, гейт первых звёзд или более позднее событие,~--- остаётся открытым (\S~\ref{sec:open_start}).

\subsection{Часы сложности}
\label{sec:complexity_clocks}

Наличие «часов сложности» следует из критического пути: при кинематически уникальной траектории (\S~\ref{sec:uniqueness}) и каскадном закреплении в наследуемой памяти (\S~\ref{sec:cascade}) накопленная сложность растёт монотонно и эмпирически экспоненциально, поэтому любая монотонная её мера~--- хронометр после калибровки темпа; «часы» здесь не дополнительный постулат.

\textbf{Первые часы: функциональный геном (Шаров).} Алексей Шаров предложил использовать рост функционального (неизбыточного) размера генома как «часы» происхождения и эволюции жизни~\cite{Sharov}: регрессия $\log_{10}G = 8{,}64 + 0{,}89\,t$ ($t$~--- млрд лет) даёт экспоненциальный рост $\sim$ в $7{,}8$ раза за 1~Gyr (удвоение $\sim$ каждые 340~млн лет), а экстраполяция к одному нуклеотиду переносит зарождение жизни на $\sim 9{,}7 \pm 2{,}5$~Gyr назад~\cite{Sharov2013}~--- ещё до образования Земли. Отсюда следует, что критический путь начинается не на планете: значительная часть усложнения~--- от пребиотической химии до первого репликатора~--- протекает в открытом космосе, на миллиарды лет раньше появления океанов.

\textbf{Независимые «часы сложности»: индекс сборки.} Теория сборки (Assembly Theory) предлагает независимую меру молекулярной сложности~--- \emph{индекс сборки} (MA), равный минимальному числу операций копирования, необходимых для построения молекулы из атомов~\cite{Marshall2021,Cronin2023}. Опубликованные значения: глицин MA$\approx$4~\cite{Marshall2021}, аденин MA$\approx$7, АТФ MA$\approx$21~\cite{Mathis2021}, типичные метаболиты MA$\approx$12--15 (аденозин MA$\approx$13~\cite{Marshall2021}). Шкала $\log_2(\mathrm{MA})$ совместима со шкалой Шарова: для случайных полимерных последовательностей MA$\approx N$ (длина цепи), поэтому $\log_2(\mathrm{MA}) \approx \log_2 N$~--- именно то, что измеряет Шаров. Это открывает возможность \textbf{унификации} двух шкал в единые «часы сложности» от Большого взрыва до наших дней: абиотическая фаза (химический рост MA) плавно переходит в биологическую (геномный рост по Шарову). Количественная оценка показывает, что обе кривые совместимы, однако скорость усложнения в абиотической фазе ($\sim 0{,}6$~Gyr$^{-1}$ в логарифмической шкале) примерно в 4--5 раз ниже, чем в биологической ($\sim 3{,}0$~Gyr$^{-1}$ по Шарову). Настоящая работа \emph{не} ставит задачи построения полной унифицированной шкалы и использует упрощённую однофазную модель суммы стадий; разработка унифицированных часов сложности вынесена в открытые вопросы (\S~\ref{sec:unified_clocks}).

\textbf{Третьи «часы»: плотность потока свободной энергии (Чейссон).} Сложность можно отслеживать не только информационно (геном, индекс сборки), но и \emph{энергетически}. Эрик Чейссон предложил \emph{плотность потока свободной энергии} $\Phi_m$ (эрг$\cdot$с$^{-1}\cdot$г$^{-1}$) как универсальную метрику сложности и движущий фактор эволюции~\cite{Chaisson2011}: она монотонно растёт вдоль всей космической лестницы: звёзды $\sim 1$, планеты $\sim 10^2$, растения $\sim 10^3$, животные $\sim 10^4$, мозг $\sim 10^5$, технологическое общество $\sim 10^6$, и, отложенная против времени появления структур, даёт примерно экспоненциальный рост на протяжении $\sim 13$~млрд лет (удвоение порядка $\sim 0{,}5$--$1$~млрд лет). В отличие от двух предыдущих мер, эти часы фиксируют не информацию, а \emph{диссипацию}, и потому напрямую сопрягаются с каскадной термодинамической картиной (\S~\ref{sec:cascade}): каждый шаг «вправо» поднимает высоту ступеньки и $\Phi_m$. \emph{Уточнение:} из механики критического пути (\S~\ref{sec:cascade}) физически растёт прежде всего \textbf{удельная скорость производства энтропии} $\dot S_{\mathrm{prod}}$, а не энергия как таковая,~--- именно $\dot S_{\mathrm{prod}}$ напрямую связана со ступеньками сложности (каждый шаг «вправо» увеличивает прирост энтропии $\Delta S$). Поток свободной энергии $\Phi_m$~--- удобная и измеримая, но по сути \emph{производная} величина; при фиксированной $T$ они пропорциональны ($\dot S = \Phi/T$), так что для роли часов это одно и то же, но физически обоснована именно скорость производства энтропии. Именно эти часы использованы ниже для энергетической оценки уровня C* из бюджета пылинки (\S~\ref{sec:idp_dust}, \S~\ref{sec:freeze}).

\textbf{К объединению трёх часов.} Три независимые меры~--- функциональный геном (Шаров), индекс сборки MA (Cronin/Walker) и плотность потока свободной энергии $\Phi_m$ (Чейссон)~--- описывают одну траекторию усложнения с разных сторон: информационной (первые две) и энергетической (третья). Их объединение в единые «часы сложности» от Большого взрыва до цивилизации выглядит естественным: все три растут примерно экспоненциально и связаны каскадной моделью (\S~\ref{sec:cascade}), где рост наследуемой информации ($N$, MA) и рост диссипации ($\Phi_m$)~--- две стороны одного подъёма по энергетическим ступенькам с закреплением результата в памяти. Согласование темпов нетривиально: информационные часы в биофазе идут в $\sim 4$--$5$ раз быстрее, чем в абиотической, а $\Phi_m$ растёт даже при почти неизменном геноме (у прокариот геном вышел на плато, тогда как диссипация продолжала расти за счёт метаболической оптимизации). Но это не противоречие: все три часов меряют \emph{одно и то же}~--- рост сложности вдоль критического пути,~--- лишь с немного разных сторон (наследуемая информация против диссипации), и потому \textbf{дополняют друг друга}. Их расхождения не обесценивают часы, а несут дополнительный смысл~--- указывают на смену ведущего носителя усложнения; вместе три меры дают более полную картину единой траектории, чем любая по отдельности. Построение согласованных трёхкомпонентных часов (с переводом между $\log N$, $\log \mathrm{MA}$ и $\log \Phi_m$ и явными пофазовыми темпами) мы оставляем отдельной задачей (\S~\ref{sec:unified_clocks}); для настоящей работы существенно лишь то, что все три независимо указывают на единый, примерно экспоненциальный и стартующий ещё до образования Земли критический путь.

\subsection{Локальная конвергентность пребиотической химии}
\label{sec:local_conv}

\textbf{Порог C*}~--- уровень преклеточной сложности, ниже которого открытая космическая эволюция устойчива, а выше которого необходимы устойчивая жидкофазная среда, ресурсный пул, концентрированный поток свободной энергии, активный трансмембранный транспорт и полноценный метаболизм~--- возможные лишь на тёплой планете. По сложности C* лежит \emph{ниже} полноценной клетки; насколько именно --- открытый вопрос (\S~\ref{sec:open_cstar_level}; для модели приняты лишь рабочие оценки, \S~\ref{sec:freeze}).

Синхронность достижения порога C* вблизи разных звёздных систем обеспечивается не обменом веществом между частями галактики~--- такой обмен работает на временных масштабах $\gg \sigma$ и физически не мог бы объяснить синхронизацию,~--- а \textbf{локальной конвергентностью} пребиотической химии. Вблизи каждой формирующейся звёздной системы пребиотические реакции протекают на одних и тех же физических входных данных:

\begin{itemize}
   \item \textbf{Химический состав:} межзвёздная среда галактического диска содержит примерно одинаковое соотношение органогенных элементов (H, C, N, O, S, P), задаваемое нуклеосинтезом предшествующих поколений звёзд. Типичный разброс ключевых отношений (прежде всего C/O) по диску для данной эпохи в солнечной окрестности~--- не более $\sim 2$--$3$ раз~\cite{Oberg2021,BrewerFischer2016}; полная металличность [Fe/H] может варьировать сильнее, но для пребиотики и ранней биологии критичны именно доступность и относительные количества органогенных элементов.
   \item \textbf{Источники энергии:} ультрафиолетовое излучение, космические лучи, ударные волны~--- универсальны и присутствуют вблизи каждой формирующейся системы.
   \item \textbf{Физические условия:} температуры молекулярных облаков и протопланетных дисков определяются одними и теми же физическими законами и схожи повсюду.
\end{itemize}

Одни и те же исходные элементы плюс одни и те же законы физики и химии порождают \textbf{конвергентный} пребиотический путь: как и в биологической эволюции (\S~\ref{sec:ltee}), самый быстрый химический путь к заданному уровню сложности воспроизводим. Каждая звёздная система \emph{локально} и \emph{независимо} «открывает» одни и те же органические молекулы в одной и той же последовательности~--- аминокислоты, нуклеооснования, рибозу, простые репликаторы~\cite{Furukawa2019,Callahan2011,Oberg2021}. Разброс в скорости определяется химической кинетикой, а не случайными особенностями конкретной системы.

\subsection{Мелкая космическая пыль}
\label{sec:idp_dust}

Где именно в космической фазе аккумулируется преклеточная сложность до C*~--- в молекулярных облаках, протопланетном диске, на кометах, в межзвёздной и межпланетной пыли (IDPs) или в метеоритах,~--- строго не установлено (\S~\ref{sec:open_carriers}). Скорее всего, накопление до C* шло на \emph{более ранних} из этих носителей~--- с большим временем жизни; межпланетная пыль у уже сформировавшейся планеты \emph{относительно молода} и живёт недолго~\cite{Nesvorny2010,Grun1985}. Для \emph{финальной} доставки готовой преклеточной химии на планету мы склоняемся к \textbf{мелкой пыли} как предпочтительному носителю: пылинок чрезвычайно много, и, в отличие от крупных метеоритов, мелкие частицы ($\sim 10$~мкм) тормозятся в атмосфере плавно и оседают на поверхность \emph{без существенного нагрева} (ниже температуры пиролиза органики $\sim 600^\circ$C), доставляя накопленную сложную химию неразрушенной; именно такая пыль вносит основную долю внеземного органического углерода на Землю~\cite{Anders1989,ChybaSagan1992,Flynn2004}.

\subsection{Заморозка космической фазы на C*}
\label{sec:freeze}
\label{sec:cstar_dust}

Конвергентная пребиотическая химия (\S~\ref{sec:local_conv}) работает до тех пор, пока шаги требуют мало энергии и не предъявляют жёстких требований к среде. Что останавливает дальнейшее усложнение выше порога C* в открытом космосе, \emph{не установлено} (\S~\ref{sec:open_freeze}); ниже~--- \emph{рабочая гипотеза} модели, справедливая для любой открытой космической площадки (пыль, кометы, метеориты и т.~д.), а не только для мелкой пыли (\S~\ref{sec:idp_dust}). Мы предполагаем, что барьер связан с \emph{условиями среды}, а не с глобальным «остыванием» Вселенной: космическое остывание слишком медленно, а энергию пребиотика на таких носителях дают локальные источники~--- УФ молодых звёзд и ударные процессы (\S~\ref{sec:local_conv}),~--- а не реликтовый фон с ранним тепловым гейтом (\S~\ref{sec:common_start}). По этой гипотезе липидная мембрана, трансмембранный транспорт и полноценный метаболизм требуют устойчивой жидкофазной среды, ресурсного пула и длительного потока свободной энергии~--- чего на открытом носителе в вакууме не обеспечить. Космическая эволюция «замораживается» на \textbf{пороговом уровне сложности} C* (преклеточные структуры): в метеоритах и протопланетных дисках находят лишь пребиотическую органику, но не живые клетки. Шаги выше C* переносятся на тёплые планеты с атмосферой, гидросферой и ресурсным пулом. Ниже~--- хронологическая и энергетическая оценки этого порога.

Экспериментальная нижняя граница клеточной сложности~--- JCVI-syn3.0~\cite{Hutchison2016}: $473$~гена, геном $\sim 5{,}31\times10^{5}$~п.н. У Шарова (\S~\ref{sec:complexity_clocks}) та же ступень «прокариот'' калибруется двумя якорями: $G \sim 5\times10^{5}$~п.н. (минимальные современные бактерии) и $t \sim 3{,}5$~Gyr по ранним ископаемым~\cite{Sharov,Sharov2013}. Мы \emph{не} отождествляем C* с этой точкой: syn3.0 задаёт лишь \emph{грубую верхнюю} оценку $G_{\text{C*}} < G_{\text{cell}} \approx 5{,}31\times10^{5}$~п.н., без независимой калибровки самого C*. На шкале Шарова C* лежит левее (меньше $G$) и, по нашему допущению, \emph{раньше} по $t$, чем первая клетка ($\sim 3{,}5$~Gyr по ископаемым). Подстановка $G_{\text{cell}}$ в регрессию $\log_{10}G = 8{,}64 + 0{,}89\,t$ не «предсказывает» C*, а лишь восстанавливает уже заданный калибровочный момент первой клетки. Между образованием Земли ($4{,}5$~Gyr) и появлением клеток ($\lesssim 3{,}5$~Gyr) календарно $\sim 1$~Gyr~--- окно планетарного участка: преклеточная химия уровня C* уже поступает с космических носителей (\S~\ref{sec:idp_dust}).

Оценим также энергобюджет минимального фронта критического пути на уровне C* по калибровке из \S~\ref{sec:resource_threshold}. Эталон~--- колба LTEE ($\sim 3 \times 10^{-4}$~Вт, $\sim 300$~мкВт), внеся две поправки. \emph{Во-первых}, эта мощность~--- с запасом: колба ($5\times10^{8}$~клеток) лежит на $\sim 2{,}7$ порядка выше порога выхода на плато скорости ($N_c\sim 10^{6}$, \S~\ref{sec:resource_threshold}), так что минимальный фронт, ещё держащийся на плато, требует лишь $\sim 10^{-6}$~Вт. \emph{Во-вторых}, на преклеточной стадии удельная диссипация на $\sim 2$ порядка ниже современной микробной (рост $\Phi_m$ по Чейссону за $\sim 4$~млрд лет, \S~\ref{sec:cascade}), что снижает фронт ещё на два порядка~--- до $\sim 10^{-9}$~Вт. Положение C* тем самым лежит примерно на $4$--$5$ порядков ниже фронта LTEE. Сравним эту мощность с энергетикой \emph{одной} пылинки (\S~\ref{sec:idp_dust}). Для зерна радиусом $r \sim 5$~мкм ($\sim 10$~мкм в диаметре) на $\sim 1$~а.е. от формирующейся звезды при потоке $F \sim 10^3$~Вт/м$^2$ перехваченная мощность $P \approx F\pi r^2 \sim 10^{-7}$~Вт; доля фотохимически активного УФ $f_{\mathrm{UV}} \sim 10^{-2}$ даёт $P_{\mathrm{UV}} \sim f_{\mathrm{UV}} P \sim 10^{-9}$~Вт~--- тот же порядок, что и минимальный фронт на уровне C*. Условие «достаточно одной пылинки» выполняется; при том же допущении о более раннем C* это согласуется с преклеточным барьером и объясняет быстрый старт жизни на Земле~--- готовая преклеточная химия была доставлена на носителях протопланетного диска (\S~\ref{sec:idp_dust}).

\subsection{Когорта земляподобных планет на критическом пути}
\label{sec:bimodal}

Наличие допланетарного участка критического пути (до порога C*, \S~\ref{sec:freeze}) определяет, какие именно планеты входят в текущий синхронизированный пик.

\textbf{Когорта критического пути} объединяет земляподобные планеты, сформировавшиеся \emph{до} появления первых клеток ($\sim 3{,}5$~млрд лет назад по ископаемым; \S~\ref{sec:freeze}). Эти планеты получили преклеточную химию уровня C* непосредственно с образованием и без разрыва запустили планетарный участок; как и весь путь, он требует локальных оазисов с избытком ресурсов (\S~\ref{sec:planetary_spread}). Стартовые условия для всех них являются common-mode: химия уровня C* в межзвёздной среде достигнута локально и независимо (\S~\ref{sec:local_conv}) вблизи каждой системы. Межпланетный разброс определяется стохастикой суммы шагов: $\sigma \approx \sqrt{\mu\,T} \sim 10^3$~лет. Это синхронизированный пик, к которому принадлежим мы.

Земля (возраст 4,5~млрд лет) находится вблизи \textbf{верхней возрастной границы когорты}: она образовалась за $\sim 1$~млрд лет до первых клеток по часам Шарова и является одной из самых «молодых» планет в этом пике. Планеты когорты должны были образоваться \emph{до} «заморозки» космической стадии на C*~--- иначе преклеточная химия не успела бы доставиться с носителей протопланетного диска (\S~\ref{sec:idp_dust}). Сформировавшиеся \emph{после} появления первых клеток в когорту не входят: планетарный участок у них стартовал с задержкой относительно синхронизированного пика.

Для оценки $n$: в текущий синхронизированный пик входят лишь земляподобные планеты старше $\sim 4$~млрд лет, что дополнительно ограничивает $n \sim 10^4$--$10^5$ (\S~\ref{sec:silence_sigma}).

%%% ===================================================================
\section{Оценка \texorpdfstring{$\sigma$}{σ} из тишины}
\label{sec:inversion}
%%% ===================================================================

\textit{Данный раздел инвертирует космическую тишину в \emph{верхние границы} на межпланетную дисперсию $\sigma$ времён выхода на детектируемую («громкую») стадию и на число $n$ планет Галактики, на которых критический путь идёт прямо сейчас (от верхней границы на $n$ следует нижняя граница на $\ell$); сопоставляет результаты с независимой астрофизикой и микробиологией и формулирует следствия для SETI.}

\subsection{Время до громкой стадии}
\label{sec:loud_time}

В согласии с разделением на «громкие» и «тихие» цивилизации~\cite{Hanson2021} мы отождествляем «громкость» с выходом на Тип~II по шкале Кардашёва ($P_{\mathrm{II}} = 10^{26}$~Вт по классической интерполяции Сагана~\cite{Kardashev1964,Sagan1973}): энергобюджет уровня звезды, дающий детектируемую на межзвёздных и межгалактических расстояниях сигнатуру (тепловое излучение сфер Дайсона и других мегаструктур~\cite{Dyson1960,Wright2014}). Это не новое предположение настоящей работы, а принятый в SETI-литературе порог детектируемости (\S~\ref{sec:open_loud}).

Текущая мощность цивилизации $P_0 \approx 2 \times 10^{13}$~Вт ($\approx 20$~ТВт на 2024 год~\cite{EI2024}), что по шкале Сагана соответствует $K \approx 0{,}73$~\cite{Sagan1973,Zhang2023}. Исторический средний темп роста мирового энергопотребления за последние десятилетия составляет $\approx 1{,}5\%$--$2{,}0\%$ в год (по данным IEA~\cite{IEA2023} и EI~\cite{EI2024}), поэтому значение $3\%$ используется здесь как реалистично-оптимистичный сценарий для ускоряющейся стадии развития техносферы; при экспоненциальном росте $3\%$ в год время до Типа~II:
\begin{equation}
t_{\mathrm{loud}} = \frac{\ln(P_{\mathrm{II}}/P_0)}{\ln(1{,}03)} \approx 990~\text{лет}.
\end{equation}
Чувствительность этой величины к темпу роста представлена в табл.~1.

\medskip
Прямых прогнозов именно на $3\%$ нет: ведущие сценарии (EIA~IEO, IEA~STEPS, ML-прогноз~\cite{Zhang2023}) дают $\sim 1$--$1{,}3\%$/год; при таком темпе выход на громкую стадию (Тип~II) занимает $\approx 2{,}3$--$2{,}9$~тыс.~лет, а при оптимистичных $3\%$~--- $\approx 990$~лет. Значение $3\%$ берётся как оптимистичный верхний край исторического опыта ($\approx 2{,}4\%$/год в среднем за 1965--2020~\cite{Smil2017,EI2024}). Текущая фаза, по-видимому, ускоряется: рост спроса дата-центров и ИИ-вычислений оценивается в $\sim 15$--$30\%$/год до 2030~\cite{IEA2025AI}, что делает сценарий ускоряющегося энергопотребления реалистичным на ближнем горизонте~--- хотя это всплеск десятилетнего масштаба в электропотреблении, а не устойчивый темп полной мощности на горизонте $t_{\mathrm{loud}}$. Существенно, что вывод от точного темпа не зависит: во всём диапазоне $1$--$4\%$ величина $t_{\mathrm{loud}}$ остаётся порядка $10^3$~лет (ср.\ табл.~1: $746$--$1964$~года при $1{,}5$--$4\%$), поэтому инвертированное $\mu = \sigma^2/T$ остаётся в окне генерационного времени микроорганизмов при любом правдоподобном темпе роста.

\begin{center}
\textbf{Таблица 1.} Время до перехода к стадии громкости ($t_{\mathrm{loud}}$) при различной скорости экспоненциального роста ($P_{\mathrm{II}} = 10^{26}$~Вт, $P_0 = 2 \times 10^{13}$~Вт)

\medskip
\begin{tabular}{c c c}
\hline
Темп роста в год & до Типа~I ($10^{16}$~Вт) & до Типа~II ($= t_{\mathrm{loud}}$) \\
\hline
1{,}5\% & 417 лет & 1964 лет \\
2{,}0\% & 314 лет & 1477 лет \\
2{,}5\% & 252 лет & 1184 лет \\
\textbf{3{,}0\%} & \textbf{210 лет} & \textbf{990 лет} \\
3{,}5\% & 181 лет & 850 лет \\
4{,}0\% & 158 лет & 746 лет \\
\hline
\end{tabular}
\end{center}

\subsection{Аргумент тишины как ограничение на дисперсию}
\label{sec:silence_sigma}

Эмпирический факт: детектируемых «громких» цивилизаций в Млечном Пути не наблюдается~--- систематические ИК-поиски частичных сфер Дайсона по данным Gaia, 2MASS и WISE (Project Hephaistos~\cite{Suazo2024}) не выявили подтверждённых мегаструктур. Статистические аргументы парадокса Ферми показывают, что отсутствие сигналов накладывает верхнюю границу на долю цивилизаций, достигших уровня, при котором они становились бы межзвёздно заметными~\cite{Wright2014,Hanson2021}. Если бы межпланетный разброс времён выхода $\sigma$ существенно превышал $t_{\mathrm{loud}}$, заметная доля близких цивилизаций опережала бы нас более чем на $t_{\mathrm{loud}}$, успела бы стать громкой и была бы видна.

Количественно примем модель с \emph{галактической зоной обитаемости} (ГЗО)~\cite{Lineweaver2001,Scherf2024}: $n$ миров равномерно в цилиндрическом кольце $r_{\min}\le r\le r_{\max}$ ($r_{\min}\approx4$~кпк~--- внутренний край: высокая частота сверхновых и избыточная металличность у балджа; $r_{\max}\approx14$~кпк~--- внешний край: дефицит металлов для скалистых планет) и полной толщины $h = 10^3$~св.~лет (объём $V_{\mathrm{GHZ}} = \pi(r_{\max}^2-r_{\min}^2)h$). Кольцо $4$--$14$~кпк~--- упрощённая ГЗО, согласуемая по порядку объёма с расчётами Scherf--Lammer~\cite{Scherf2024}; при существенном сужении кольца оценка тишины ужесточается, не меняя порядковых выводов ($n\sim10^4$--$10^5$, $\sigma\sim10^3$~лет). Сглаженный профиль плотности $\propto e^{-r/R_d}$ ($R_d=3$~кпк) вместо резких границ кольца смещает $\sigma$ на $\sim1$--$2\%$ (Monte Carlo, $N_{\mathrm{MC}}=4\times10^5$). Наблюдатель~--- на солнечном радиусе $R_\odot=8$~кпк, $c = 1$~св.~год/год. Подынтегральный множитель $\Phi(-(t_{\mathrm{loud}}+s)/\sigma)$ по расстоянию $s$ спадает на масштабе нескольких $\sigma_r$, где $\sigma_r\equiv c\sigma\sim10^3$~св.~лет ($\sim0{,}3$~кпк; $\sigma$~--- дисперсия \emph{календарных} времён выхода на Тип~II между планетами, $s$ и $\sigma_r$~--- в св.~годах при $c=1$~св.~год/год), а $R_\odot$ удалён от обоих краёв кольца на $\gg\sigma_r$; в эффективной окрестности наблюдателя плотность \emph{постоянна}. Мир на трёхмерном расстоянии $s$ от наблюдателя учитывается в интеграле как \emph{потенциально} видимый сегодня, если опередил нас не меньше чем на $t_{\mathrm{loud}} + s$~--- построил громкую фазу ($t_{\mathrm{loud}}$) плюс передал сигнал ($s/c$),~--- с вероятностью $\Phi(-(t_{\mathrm{loud}}+s)/\sigma)$ (верхний хвост распределения времён выхода на Тип~II; следствие ЦПТ). Ожидаемое число таких соседей:
\begin{equation}
\boxed{\;\mathbb{E}[V] = \frac{n}{V_{\mathrm{GHZ}}}\int_{-h/2}^{h/2}\!\!\int_{r_{\min}}^{r_{\max}} \Phi\!\left(-\frac{t_{\mathrm{loud}}+s(r,z)}{\sigma}\right) 2\pi r\,\mathrm{d}r\,\mathrm{d}z \;\lesssim\; 1,\;}
\end{equation}
где $s(r,z)$~--- расстояние от $(R_\odot,0,0)$ до точки $(r,z)$. При расстоянии от $R_\odot$ до краёв кольца $\gg\sigma_r$ интеграл эквивалентен однородной плотности $\rho=n/V_{\mathrm{GHZ}}$ и берётся в замкнутом виде: с $a=t_{\mathrm{loud}}/\sigma$ и $\beta=h/2\sigma_r$ переход к 3D-расстоянию $s$ (при $R\to\infty$, поправка $\sim e^{-(R/\sigma_r)^2/2}$ пренебрежима) даёт $\mathbb{E}[V]=(n/V_{\mathrm{GHZ}})\,4\pi\sigma_r^3\,W(a,\beta)$, $W(a,\beta)=\int_0^\beta K(a;y)\,dy$, где $K(a;y)=\int_y^\infty x\Phi(-(a+x))dx=\tfrac12[(1+a^2-y^2)\Phi(-(a+y))+(y-a)\varphi(a+y)]$; внешний интеграл~--- конечная комбинация $\Phi$ и $\varphi$ в точках $a$ и $a+\beta$ (гауссовы моменты). Предельные случаи: $\beta\to0$ даёт 2D-форму $\sigma_r^2 K(a)$, $K(a)=\tfrac12[(1+a^2)\Phi(-a)-a\varphi(a)]$; $\beta\to\infty$~--- 3D-форму $\sigma_r^3 M(a)$, $M(a)=\tfrac13[(a^2+2)\varphi(a)-a(3+a^2)\Phi(-a)]$. Конечная толщина $h$ существенна: на малых расстояниях $s \lesssim h/2$ соседи распределены трёхмерно, на больших~--- планарно. Из $\mathbb{E}[V]\lesssim 1$ для каждого $n$ находят \textbf{наибольшее} $\sigma$, при котором тишина ещё возможна (верхняя граница, табл.~2); $n$~--- число планет в ГЗО Галактики, на которых критический путь идёт прямо сейчас, а $\sigma$~--- межпланетный разброс времён \emph{выхода на Тип~II}, а не момент, когда мы их \emph{увидим}.

Важно явно зафиксировать тип ограничения. Тишина задаёт \emph{верхние} границы на $\sigma$ и $n$; ни при слишком малой $\sigma$ (все выходят почти синхронно), ни при $n=0$ (мы одни) молчание не нарушается. Количественно: для каждого $n$ табл.~2 даёт наибольшее допустимое $\sigma$; при фиксированном $\sigma$ тишина \emph{исключает} $n$ на порядки выше $\sim10^5$. Нижних границ на $\sigma$ и $n$ нет. Поскольку $\ell \propto n^{-1/3}$ (табл.~2), от верхней границы на $n$ следует \emph{нижняя} граница на $\ell$ ($\gtrsim 220$~св.~лет при $n \lesssim 10^5$); верхней границы на $\ell$ тишина не даёт. Оценка $\mu=\sigma^2/T$ (\S~\ref{sec:mu_from_sigma})~--- следствие верхней границы на $\sigma$ (сводка~--- табл.~4).

Именно здесь~--- и только здесь~--- модель использует \emph{форму} распределения (гауссов $\Phi$), а не одну лишь концентрацию. Это самостоятельное допущение, причём задействуется оно в \emph{верхнем хвосте}, где сходимость суммы к нормали по центральной предельной теореме наиболее медленная: реальные крупные уклонения могут отклоняться от гауссовых. Внутри модели поправку к гауссовому хвосту можно оценить разложением Эджворта: асимметрия суммы $\gamma_\Sigma=\gamma_{\text{шаг}}/\sqrt{N_{\mathrm{eff}}}$, $N_{\mathrm{eff}}=(\sum_i\sigma_i^2)^2/\sum_i\sigma_i^4$; при $\gamma_{\text{шаг}}\sim1$ (экспоненциальный шаг, $c_v\approx1$, \S~\ref{sec:mu_from_sigma}) и $N_{\mathrm{eff}}\sim N\sim10^{14}$ относительный сдвиг границы населённости $\Delta n/n\sim1/\sqrt{N_{\mathrm{eff}}}\sim10^{-7}$ в рабочей зоне интеграла ($z\approx1{,}5$--$2$). Заметное отклонение потребовало бы $N_{\mathrm{eff}}\lesssim10^2$, т.е. скрытого доминирующего шага (\S~\ref{sec:heterogeneity}, \S~\ref{sec:planetary_steps})~--- того же класса, что проверяет тест эукариогенеза (\S~\ref{sec:eukaryo}). На саму $\sigma$ хвостовая форма влияет слабо: фактор-$10$ по $n$ соответствует лишь $\sim1{,}7$-кратному сдвигу $\sigma$.

\begin{center}
\textbf{Таблица 2.} Верхняя граница $\sigma$ при $\mathbb{E}[V]=1$ (ГЗО: $r_{\min}=4$~кпк, $r_{\max}=14$~кпк, $h = 10^3$~св.~лет, $R_\odot=8$~кпк, $t_{\mathrm{loud}} = 990$~св.~лет). $n$~--- миры равномерно в объёме $V_{\mathrm{GHZ}}$; $\sigma$~--- решение замкнутой формы выше. $\ell = \Gamma(4/3)\,(3(r_{\max}^2-r_{\min}^2)h/(4n))^{1/3}$~--- среднее 3D-расстояние до ближайшего соседа ($\rho = n/V_{\mathrm{GHZ}}$; $\ell \propto n^{-1/3}$, от верхней границы на $n$~--- нижняя на $\ell$, табл.~4).

\medskip
\small
\begin{tabular}{c c c}
\hline
$n$ & $\sigma$, лет & $\ell$, св. лет \\
\hline
$10^2$ & 6960 & 2170 \\
$10^3$ & 2690 & 1010 \\
$10^4$ & 1290 & 470 \\
\textbf{$2{,}9 \times 10^4$} & \textbf{990} & \textbf{330} \\
$10^5$ & 775 & 220 \\
$10^6$ & 550 & 101 \\
$10^7$ & 430 & 50 \\
$10^8$ & 350 & 22 \\
\hline
\end{tabular}
\normalsize
\end{center}

\subsection{Число планет на критическом пути}
\label{sec:kepler}

\textbf{Оценка Scherf--Lammer 2024.} Работа~\cite{Scherf2024,Lammer2024} оценивает максимальное число \emph{землеподобных местообитаний} (Earth-like Habitats)~--- скалистых планет в зоне обитаемости сложной жизни, способных удержать устойчивую азот-кислородную атмосферу земного типа с малым $\mathrm{CO_2}$,~--- с учётом истории звездообразования, начальной функции масс, галактической зоны обитаемости и термической устойчивости атмосфер. Результат: во всём Млечном Пути таких сред \textbf{не более $n_{\text{EH}} \sim 6\times10^4$--$2{,}5\times10^5$}, преимущественно у G- и K-звёзд. Эта оценка уже включает ключевые срезы~--- возраст планеты, поколение и металличность звезды~\cite{Lineweaver2001,Behroozi2015}, положение в галактической зоне обитаемости и стабильность атмосферы. Полный набор критериев отбора: скорость звездообразования и начальная функция масс, галактическая зона обитаемости, преимущественно G- и K-звёзды, зона обитаемости \emph{сложной} жизни, частота скалистых планет ($\eta_\oplus$), одновременное наличие океанов и суши, вероятное требование крупной луны и термическая устойчивость N$_2$--O$_2$-атмосферы при заданной доле CO$_2$.

\textbf{Совпадение с тишиной.} Независимый результат из космической тишины (\S~\ref{sec:silence_sigma}, табл.~2)~--- $n \sim 10^4$--$10^5$ цивилизационных планет при $\sigma \sim 0{,}7$--$1{,}3$~тыс.~лет~--- совпадает с астрофизической оценкой Scherf--Lammer ($n_{\text{EH}} \sim 10^5$). Два независимых канала~--- астрофизика населённости и статистика молчания~--- сходятся на $n \sim 10^5$; инверсия тишины (\S~\ref{sec:silence_sigma}) и Scherf--Lammer считают $n$ в одном классе объёмов~--- кольце галактической зоны обитаемости. Будь планет на критическом пути на порядки больше, тишина была бы нарушена (\S~\ref{sec:silence_sigma}). Поскольку $n_{\text{EH}}$~--- верхняя оценка числа пригодных сред, а не гарантированных цивилизаций, фактическое $n$ на критическом пути может быть ниже; это только ужесточает согласие с тишиной.

\subsection{Средний элементарный такт \texorpdfstring{$\mu$}{μ}}
\label{sec:mu_from_sigma}

При однофазном приближении с единым усреднённым тактом ($\sigma = \sqrt{\mu T}$) верхняя граница на $\sigma$ (из условия тишины) \textbf{инвертируется} в оценку \emph{среднего} времени элементарного шага $\mu \equiv T/N$ (\S~\ref{sec:concentration}), а не длительности каждого отдельного шага ($\mu_i$ могут различаться):
\begin{equation}
\boxed{\mu = \frac{\sigma^2}{T}.}
\end{equation}
При $T = 1{,}38 \times 10^{10}$~лет для каждой строки табл.~2 (тишина при $\mathbb{E}[V]=1$) получаем согласованные параметры (табл.~3): $\mu=\sigma^2/T$, $N=(T/\sigma)^2$ и темп роста как \emph{обратную} подстановку $t_{\mathrm{loud}}=\sigma$ по формуле табл.~1 (не вход интеграла, где $t_{\mathrm{loud}}=990$~лет). Отождествление $\sigma \approx t_{\mathrm{loud}}$ и сценарий $3\%$/год выполняются лишь в выделенной строке.

Объединяя концентрацию $\sigma = \sqrt{\mu T}$ с условием тишины на пределе видимости $\sigma \approx t_{\mathrm{loud}}$ (\S~\ref{sec:silence_sigma}), получаем замкнутое выражение
\begin{equation}
\mu \approx \frac{t_{\mathrm{loud}}^{2}}{T},
\end{equation}
в которое входят только \emph{небиологические} величины: горизонт выхода на громкую фазу $t_{\mathrm{loud}} \sim 10^3$~лет (наша собственная энергетическая траектория, \S~\ref{sec:loud_time}) и возраст Вселенной $T$. Подстановка $t_{\mathrm{loud}} \approx 10^3$~лет и $T \approx 1{,}38 \times 10^{10}$~лет даёт $\mu \sim 40$~мин. Без постулата $c_v = 1$ общее соотношение $\sigma^2 = c_v^2\mu T$ (\S~\ref{sec:cascade}) даёт $\mu = \sigma^2/(c_v^2 T) = (40~\text{мин})/c_v^2$ и $N = c_v^2 (T/\sigma)^2$. Однако $c_v$ не обязан задаваться извне: соотношение $\sigma^2 = c_v^2\mu T$ связывает три величины, и любые две задают третью. Подставив \emph{независимо} известные $\mu \approx 40$~мин (время деления, микробиология) и $\sigma \approx 990$~лет (тишина, астрономия), получаем $c_v^2 = \sigma^2/(\mu T) \approx 0{,}93$, то есть $c_v \approx 0{,}97 \approx 1$. Кинетически это согласуется с тем, что шаг лимитируется единственным переходом с постоянной интенсивностью: минимум параллельных попыток экспоненциален, $c_v=1$ точно (речь о статистике \emph{ожидания} шага, а не об удержании результата, \S~\ref{sec:ratchet}). Иными словами, при двух \emph{независимо} заданных входах ($\sigma$ из тишины, $\mu$ из микробиологии) третья величина $c_v$ оказывается близкой к единице, а экспоненциальная форма одношагового ожидания не постулируется произвольно, а \emph{согласуется} с этой парой каналов (астрономия $\leftrightarrow$ микробиология) и кинетикой §~\ref{sec:ratchet}. Здесь речь о \emph{форме} распределения шага ($c_v$), а не о выборе опорной строки табл.~3~--- это отдельный, более слабый шаг (ниже). Сама концентрация $\sigma/T = c_v/\sqrt{N}$ от значения $c_v$ не зависит.

\begin{center}
\textbf{Таблица 3.} Сводная оценка по строкам табл.~2 ($T = 1{,}38 \times 10^{10}$~лет, $\mu = \sigma^2/T$~--- среднее время шага). $n$~--- число планет на критическом пути в ГЗО Галактики (\S~\ref{sec:silence_sigma}); $N$~--- число элементарных шагов; темп роста~--- скорость экспоненциального роста $P(t)$, при которой $t_{\mathrm{loud}}=\sigma$ (обратная подстановка из табл.~1, не вход интеграла тишины); самосогласована выделенная строка ($\sigma=t_{\mathrm{loud}}=990$~лет, $3\%$/год).

\medskip
\small
\begin{tabular}{c c c c c c}
\hline
$n$ & $\sigma$, лет & $\ell$, св.~лет & рост, \%/год & $\mu$, мин (ср.) & $N$ \\
\hline
$10^2$ & 6960 & 2170 & 0{,}42 & 1840 & $3{,}9 \times 10^{12}$ \\
$10^3$ & 2690 & 1010 & 1{,}1 & 276 & $2{,}6 \times 10^{13}$ \\
$10^4$ & 1290 & 470 & 2{,}3 & 63 & $1{,}1 \times 10^{14}$ \\
\textbf{$2{,}9 \times 10^4$} & \textbf{990} & \textbf{330} & \textbf{3{,}0} & \textbf{37} & \textbf{$1{,}9 \times 10^{14}$} \\
$10^5$ & 775 & 220 & 3{,}8 & 23 & $3{,}2 \times 10^{14}$ \\
$10^6$ & 550 & 101 & 5{,}3 & 12 & $6{,}3 \times 10^{14}$ \\
$10^7$ & 430 & 50 & 6{,}8 & 7{,}0 & $1{,}0 \times 10^{15}$ \\
$10^8$ & 350 & 22 & 8{,}4 & 4{,}7 & $1{,}6 \times 10^{15}$ \\
\hline
\end{tabular}
\normalsize
\end{center}

Полученная инверсией оценка $\mu \approx 40$~мин (в опорной точке табл.~3~--- $\approx 37$~мин; величина порядковая, интервал правдоподобия $25$--$75$~мин~--- \S~\ref{sec:intervals}) лежит в диапазоне генерационного времени микроорганизмов (\emph{E.~coli} $\sim 20$--$60$~мин): это такт \emph{плато скорости} микробной фазы (\S~\ref{sec:resource_threshold}), а не календарное поколение многоклеточных (\S~\ref{sec:sex}). Космическая тишина (астрономия) и время деления клетки (микробиология) согласуются по \emph{порядку величины} с одной и той же «единицей такта» эволюции. При ином темпе роста (столбец «рост» в табл.~3) значение $\mu$ вышло бы из биологически осмысленного окна. Единый усреднённый $\mu$~--- упрощение однофазной инверсии: $\mu=\sigma^2/T$ даёт путевое $\mu^2$-среднее ($\sigma^2=\sum_i N_i\mu_i^2$), а не такт отдельной фазы; часы Шарова~\cite{Sharov2013} и индекс сборки~\cite{Marshall2021,Cronin2023} указывают на $\geq 2$ режима роста сложности (абиотический $\sim 0{,}6$~Gyr$^{-1}$, биологический $\sim 3{,}0$~Gyr$^{-1}$). При $N = T/\mu \approx 1{,}9\times10^{14}$ (табл.~3) по времени $\sim 65\%$ шагов~--- пребиотика ($\sim 9$~Gyr), ранняя планетная $\sim 9\%$, микробная $\sim 22\%$ ($\sim 3$~Gyr одноклеточной фазы на Земле~\cite{Nutman2016}), многоклеточная и цивилизационная $\sim 4\%$; если собственный такт пребиотики отличается от биологического в $\sim 4$--$5$ раз, инвертированное $\mu$ смещается на фактор \emph{в несколько раз}, тогда как $\sigma$ из тишины не затрагивается~--- ослабевает лишь сопоставление $\mu$ с микробным окном (темп пребиотических шагов~--- \S~\ref{sec:open_prebiotic_tempo}).

\textbf{Оговорка о выборе строки.} Инверсия «тишина $\to \sigma$» при фиксированном $T$ задаёт $\mu=\sigma^2/T$ \emph{для каждой} строки табл.~2 (табл.~3), а не единственное $\mu$ из первых принципов. Из всех строк биологически осмысленна лишь узкая~--- $\mu \approx 20$--$60$~мин, отвечающая $n \sim 10^4$--$10^5$ и темпу роста энергопотребления $\sim 2$--$3\%$/год. Опорную строку ($n \approx 2{,}9\times10^4$, рост $3\%$) мы выбираем не как единственно возможную, а как \emph{наиболее правдоподобную}: темп $3\%$~--- реалистично-оптимистичный сценарий для ускоряющейся техносферы, а полученное $\mu \approx 37$~мин совпадает с измеренным временем деления микроорганизмов. Совпадение \emph{отсекает} строки с невероятным тактом и \emph{отбирает} ту, где сходятся три эмпирических ориентира~--- микробиология ($\mu$), астрофизическая оценка населённости (Scherf--Lammer, $n$) и реалистичный темп роста энергопотребления; это усиливает оценку, но не доказывает её. Через $\sigma = \sqrt{\mu T}$ то же окно $\mu \approx 20$--$60$~мин отвечает $\sigma \approx 730$--$1260$~лет; по табл.~2 такой диапазон $\sigma$ допускает при наблюдаемой тишине $n \sim 10^4$--$10^5$ цивилизационных планет в Млечном Пути. Это совместное фальсифицируемое предсказание двух независимых каналов: если населённость Галактики окажется $n \gg 10^5$, то либо $\sigma < 730$~лет (тогда $\mu$ быстрее времени деления клетки~--- физически сомнительно), либо тишина обязана нарушиться.

\subsection{Интервалы неопределённости и устойчивость $\sigma$}
\label{sec:intervals}

Приводимый ниже интервал~--- \emph{не} частотный доверительный интервал (здесь нет выборочной статистики), а \emph{диапазон правдоподобия}, полученный распространением неопределённости входных параметров через соотношения $\mu=\sigma^2/(c_v^2T)$ и $N=c_v^2(T/\sigma)^2$. Он \emph{условен}: задаётся при $n\in[10^4,10^5]$ (пересечение астрофизической оценки Scherf--Lammer, \S~\ref{sec:kepler}, и \emph{верхней} границы тишиной на $n$, табл.~2) и $c_v\approx1$ (\S~\ref{sec:mu_from_sigma}). Столбец «диапазон» в табл.~4~--- окно правдоподобия внутри этих ограничений, а не симметричный доверительный интервал; типы границ тишины ($\le$ на $\sigma$ и $n$, $\ge$ на $\ell$) указаны в последнем столбце и в заголовке табл.~4.

\textbf{Слабая зависимость $\sigma$ от $n$~--- ключевое свойство.} Верхняя граница $\sigma(n)$ из тишины (табл.~2) \emph{сильно сублинейна}: при изменении населённости в \textbf{10 раз} (от $10^4$ до $10^5$) $\sigma$ меняется лишь в $\approx1{,}7$ раза (с $1290$ до $775$~лет), тогда как $\mu$ и $N$~--- в $\approx2{,}8$ раза (как $\sigma^2$). Поэтому даже грубая неопределённость $n$ \emph{не} расшатывает верхнюю границу на $\sigma$: дисперсия времён выхода зафиксирована на уровне $\sim10^3$~лет почти независимо от того, сколько именно планет идёт по критическому пути. Это же делает предсказание $\mu\sim40$~мин устойчивым~--- оно остаётся в окне деления микроорганизмов ($20$--$60$~мин) при любой правдоподобной $n$.

\textbf{Слабая зависимость $\ell$ от $n$.} Поскольку $\ell \propto n^{-1/3}$, в том же окне $n\in[10^4,10^5]$ среднее расстояние до ближайшего соседа меняется лишь в $\approx1{,}4$ раза ($470$--$220$~св.~лет). Единственное ограничение тишины на $\ell$~--- \emph{нижняя} граница ($\gtrsim 220$~св.~лет при $n\lesssim 10^5$); при меньшем $n$ соседи могут быть и дальше. Опорное $\ell\sim330$~св.~лет (табл.~2) согласуется с интерпретацией SETI (\S~\ref{sec:seti}): близкие соседи в когорте есть, но ещё не «громкие».

\begin{center}
\textbf{Таблица 4.} Интервалы правдоподобия при $c_v\approx1$, $T=1{,}38\times10^{10}$~лет. Тишина (\S~\ref{sec:silence_sigma}) задаёт \emph{верхние} границы на $\sigma$ и $n$ (нижних границ на них нет); от верхней границы на $n$ следует \emph{нижняя} граница на $\ell$; $\mu$ и $N$~--- следствие верхней границы на $\sigma$.

\medskip
\small
\begin{tabular}{l c c c}
\hline
Параметр & Центр & Диапазон & Главный источник \\
\hline
$\sigma$ & $1000$~лет & $800$--$1400$~лет & тишина ($\le$), $n$ \\
$n$ & $3\times10^4$ & $10^4$--$10^5$ & тишина ($\le$), Scherf--Lammer \\
$\mu$ & $40$~мин & $25$--$75$~мин & $\sigma$ ($\sigma^2$) \\
$N$ & $1{,}9\times10^{14}$ & $(1{,}0$--$2{,}9)\times10^{14}$ & $\sigma$ ($\sigma^{-2}$) \\
$\ell$ & $330$~св.~лет & $220$--$470$~св.~лет & $n$ ($\ge$) \\
\hline
\end{tabular}
\normalsize
\end{center}

\subsection{Следствие для SETI}
\label{sec:seti}

Статистическая синхронизация меняет постановку вопроса. Обычно молчание Ферми трактуют в одной из двух крайностей: либо мы одни (или крайне редки), либо другие цивилизации значительно старше и обогнали нас. Модель предлагает третий ответ: \textbf{большинство цивилизаций возникают примерно одновременно}, поскольку концентрация суммы большого числа шагов сжимает разброс до $\sigma \sim 10^3$~лет~--- пренебрежимо мало на масштабе жизни Вселенной, так что ни одна цивилизация не оказывается значимо «старше» прочих. Отсюда другая интерпретация SETI: другие технологические цивилизации, если они есть, с высокой вероятностью находятся на \emph{сопоставимой} стадии~--- то есть являются \textbf{потенциальными конкурентами}, а не маяками; молчание объясняется не отсутствием или удалённостью других, а тем, что все они появились недавно (по космическим меркам) и ещё не вышли на детектируемый уровень. Стратегии SETI, ориентированные на поиск сигналов значительно более развитых цивилизаций, могут оказаться принципиально неверными.

%%% ===================================================================
\section{Предсказания и фальсифицируемость}
\label{sec:predictions}
%%% ===================================================================

\textit{Данный раздел формулирует проверяемые предсказания модели по нескольким независимым фронтам; каждое вынесено в отдельный подраздел с явным условием фальсификации.}

\subsection{Эукариогенез (решающий тест)}
\label{sec:eukaryo}

Гипотеза, что эукариотическая клетка на Земле возникла однократно за $\sim 2$~млрд лет, выдвигает эукариогенез как главного кандидата на уникальный медленный «трудный шаг» и самый острый тест модели. Контртезис: это не скачок, а длинная синтрофная цепочка высоковероятных подшагов~--- асгардские археи с эукариотическими сигнатурами~\cite{Spang2015,Zaremba2017}, культивируемый \emph{Ca.}~Prometheoarchaeum syntrophicum (модель E3: Entangle--Engulf--Endogenize)~\cite{Imachi2020}, синтрофная и водородная гипотезы~\cite{LopezGarcia2020,Martin1998}, актиновые гомологи у Lokiarchaeota~\cite{Spang2015}, наблюдаемый актиновый цитоскелет у Asgard-археи~\cite{Rodrigues2023}, доминантный вклад Asgard-архей и пошаговый переход FECA$\to$LECA~\cite{Tobiasson2026,Speijer2025,BravoArevalo2025}. Тогда длительность $\sim 2$~млрд лет~--- длина цепочки, а не один трудный шаг; полной механистической модели пока нет, поэтому тест остаётся \emph{открытым}. \emph{Фальсификация:} если хотя бы одно звено окажется устойчиво неразложимым, уникальным и медленным~--- не проходящим быстрее $\sim 10^3$~лет даже с учётом всей популяции параллельно эволюционирующих линий (параллельный поиск его не ускоряет),~--- это даёт большой разброс времён выхода и разрушает синхронность.

\subsection{Возраст планеты и стадия жизни}
\label{sec:pred_planet_age}

Достигнутая стадия эволюции жизни монотонно растёт со временем, проведённым на критическом пути, поэтому коррелирует с возрастом планеты. На планетах \emph{старше} Земли (когорта, стартовавшая без разрыва, \S~\ref{sec:bimodal}), при избытке ресурсов и условиях, достаточных для критического пути (\S~\ref{sec:resource_threshold}, \S~\ref{sec:plateau}; современный Марс сюда не входит), жизнь, если она есть, уже не может находиться на простой одноклеточной стадии: за время не меньшее земного она прошла её и достигла сложных, технологических форм. В силу синхронности (\S~\ref{sec:seti}) такие планеты не на эпохи впереди, а в пределах $\sigma \sim 10^3$~лет от нашего уровня. Напротив, простая одноклеточная жизнь ожидается \emph{только} на планетах, существенно \emph{моложе} Земли~--- сформировавшихся уже после достижения порога C* и потому отставших от синхронизированного пика (\S~\ref{sec:bimodal}), которые сейчас находятся в ранней фазе пути. \emph{Фальсификация:} обнаружение богатой \emph{простой} (микробной, доцивилизационной) биосферы на землеподобной планете заметно старше Земли при избытке ресурсов и тех же необходимых условиях~--- без признаков продвижения к сложности~--- противоречило бы монотонности и синхронности критического пути.

\subsection{Конвергентная форма цивилизации}
\label{sec:pred_convergent_form}

Критический путь задаёт самую быструю траекторию через пространство сложности \emph{по полному времени} (\S~\ref{sec:uniqueness}); вместе с конвергентной эволюцией (\S~\ref{sec:ltee}) это влечёт как \emph{логическое следствие} не только синхронность времён, но и \emph{частичную} сходимость \emph{формы цивилизации} (\S~\ref{sec:time_form}): технологические цивилизации на других планетах \emph{ожидаемо} сходятся к воспроизводимому \emph{функциональному} плану биологического носителя~--- манипуляторные конечности, крупный социальный мозг, орудийность,~--- то есть к гуманоидам нашего типа, а не к рептильно-динозавровым или принципиально иным сценариям. \emph{Объём утверждения.} Конвергентна лишь ведущая линия к технологическому разуму (\S~\ref{sec:uniqueness}), а не вся биосфера: сопутствующие флора и фауна на других планетах могут быть произвольно иными~--- предсказание касается только биологического носителя цивилизации и его предковой линии, но не состава экосистемы. \emph{Подтверждение:} обнаружение на другой планете технологической цивилизации с носителем функционально гуманоидного типа~--- и при этом \emph{никого} существенно более развитого~--- подтвердило бы конвергенцию формы цивилизации и синхронность времён ($\sigma \sim 10^3$~лет, когда ни одна цивилизация не ушла вперёд на космологические сроки). Именно совмещение этих двух условий~--- «та же функциональная форма цивилизации» и «никто не опередил»~--- отличает наше предсказание от сценариев, где иные цивилизации либо принципиально иной формы, либо значительно старше и развитее нас (\S~\ref{sec:seti}). \emph{Фальсификация:} обнаружение технологической цивилизации на принципиально ином эволюционно-морфологическом базисе; обнаружение цивилизации, существенно более развитой, чем наша, разрушило бы тезис синхронности.

\subsection{Подъём по Кардашёву к Типу~II (наше будущее)}
\label{sec:pred_kardashev}

Выход на «громкую» стадию идёт через устойчивый экспоненциальный рост энергопотребления (исторически $\sim 1{,}5$--$3\%$ в год) и достигает Типа~II за $t_{\mathrm{loud}} \sim 10^3$~лет (\S~\ref{sec:loud_time}). Модель трактует это как \emph{типовой} исход критического пути и тем самым предсказывает наше собственное будущее~--- продолжение подъёма Тип~I\,$\to$\,II на горизонте $\sim 10^3$~лет (одновременно посылка о «громкости» цивилизаций и проверяемое во времени следствие; форвард-прогноз~--- \S~\ref{sec:open_loud}; галактическая конкуренция после Типа~II~--- \S~\ref{sec:galactic_competition}). \emph{Фальсификация:} устойчивое отклонение траектории от такого роста~--- долгое плато, спад, «великий фильтр» ниже Типа~II у нас и у наблюдаемых цивилизаций~--- лишает силы $t_{\mathrm{loud}}$, инверсию «тишина $\to \sigma$» и саму посылку аргумента тишины.

\subsection{Детектируемость Типа~I}\label{sec:pred_type1}

Аргумент тишины (\S~\ref{sec:silence_sigma}) калиброван на пороге Типа~II ($P_{\mathrm{II}} = 10^{26}$~Вт, $t_{\mathrm{loud}} \approx 990$~лет при $3\%$/год); при нём $\mathbb{E}[V] \lesssim 1$, и наблюдаемый ноль громких сигналов правдоподобен. Тот же интеграл с заменой $t_{\mathrm{loud}}$ на время до Типа~I ($P_{\mathrm{I}} = 10^{16}$~Вт; табл.~1: $t_{\mathrm{I}} \approx 210$~лет при $3\%$/год) и тем же $\sigma$ из табл.~2 (верхняя граница тишины, табл.~4) даёт \emph{условное} предсказание: если построить инструмент, способный межзвёздно детектировать цивилизации уровня Типа~I (техносфера с энергобюджетом планеты), то при $n \sim 10^4$--$10^5$ (сходящаяся оценка модели: тишина + Scherf--Lammer, \S~\ref{sec:kepler}) и опорных параметрах ($n \approx 2{,}9\times10^4$, $\sigma \approx 990$~лет) ожидаемое число уже видимых соседей~--- $\mathbb{E}[V_{\mathrm{I}}] \approx 5$, диапазон $\approx 3$--$9$; вероятность обнаружить \emph{хотя бы одного}~--- $\approx 1 - e^{-\mathbb{E}[V_{\mathrm{I}}]} \gtrsim 99\%$. Условие видимости то же: опережение не меньше $t_{\mathrm{I}} + s$ (выход на Тип~I плюс ход света). Среднее расстояние до ближайшего соседа $\ell \sim 330$~св.~лет ($\gtrsim 220$~св.~лет~--- нижняя граница из верхней границы на $n$, \S~\ref{sec:intervals}) согласуется с этой картиной: соседи в когорте есть, но ещё не «громкие» на пороге Типа~II. Это \emph{не} утверждение, что Тип~I уже межзвёздно детектируем при нынешних представлениях SETI, а следствие синхронности когорты: доля соседей, опередивших нас на $t_{\mathrm{I}}$, велика, но лишь близкие попадают в световой конус; при пороге Тип~II таких соседей $\lesssim 1$. \emph{Отсутствие сигналов не фальсифицирует модель:} $n$ из тишины~--- лишь верхняя граница, и при малом $n$ (вплоть до $n=0$) нулевой результат совместим с наблюдаемой тишиной по Типу~II. \emph{Подтверждение} (при подтверждённой чувствительности к Типу~I): обнаружение хотя бы одной цивилизации Типа~I в ГЗО с параметрами, согласующимися с синхронной когортой ($\sigma \sim 10^3$~лет, $n \sim 10^4$--$10^5$),~--- прямое следствие предсказанного числа $\mathbb{E}[V_{\mathrm{I}}] \approx 5$.

%%% ===================================================================
\section{Открытые вопросы}
\label{sec:open}
%%% ===================================================================

\textit{Данный раздел систематизирует нагруженные допущения, упрощения и нерешённые вопросы модели; каждый подраздел вынесен с явной отсылкой из основного текста.}

\subsection{Субгенерационный такт в многоклеточной фазе}
\label{sec:macro_tempo}

Сведение вклада многоклеточной фазы к $\sigma \sim 10^3$~лет через $\mu_{\mathrm{eff}}$ и субгенерационные каналы~--- \S~\ref{sec:sex}. Открытый вопрос~--- количественная доля субгенерационных каналов (обучение, культура, эффект Болдуина) в фанерозойной сложности; проверка~--- темпы конвергентной макроэволюции и роль пластичности в направлении генетической фиксации.

\subsection{Унифицированные часы сложности}
\label{sec:unified_clocks}

Функциональный геном (Шарова), индекс сборки MA и плотность потока свободной энергии $\Phi_m$ (Чейссон) описывают один критический путь с информационной и энергетической сторон (\S~\ref{sec:complexity_clocks}), но согласование их темпов по фазам (абиотика, биология, техносфера) нетривиально. Перспективная задача~--- построить согласованные трёхкомпонентные «часы сложности» от Большого взрыва до цивилизации: явный перевод между $\log N$, $\log \mathrm{MA}$ и $\log \Phi_m$, калибровка пофазных темпов и проверка, насколько такая шкала улучшает однофазную модель суммы стадий с единым усреднённым тактом (\S~\ref{sec:mu_from_sigma}). Настоящая работа сознательно использует упрощение; полная унификация~--- задел на будущие исследования.

\subsection{Старт критического пути}
\label{sec:open_start}

В работе за полное время $T$ принят отсчёт от теплового гейта CMB (\S~\ref{sec:common_start}, $\sim 10$--$17$~млн лет~\cite{Loeb2014}); альтернативой служит материальный гейт первых звёзд~\cite{Naoz2006} ($\sim 30$~млн лет после Большого взрыва для первой наблюдаемой звезды). Открытым остаётся, какое из этих событий~--- или более позднее~--- следует считать истинным общим стартом критического пути.

\subsection{Скорость пребиотических шагов}
\label{sec:open_prebiotic_tempo}

Лабораторные репликаторы (шаблонная лигация фон Кидровского, пептидные репликаторы, системная химия) реплицируются преимущественно \emph{параболически}, а не экспоненциально, из-за ингибирования продуктом~\cite{Adamski2020}~--- что согласуется с тем, что добиотические шаги медленнее и шумнее; но чистого универсального времени шага для этой стадии данные пока не дают.

\subsection{Уровень сложности C*}
\label{sec:open_cstar_level}

Какой именно уровень сложности обозначает C*? В \S~\ref{sec:local_conv} порог задан \emph{функционально}~--- ниже него открытая космическая эволюция устойчива, выше требуются планетные условия; установлено лишь, что C* \emph{ниже} полноценной клетки, но не \emph{насколько}. Независимой калибровки по геному ($G$), индексу сборки (MA), $\Phi_m$ или числу элементарных шагов до первой клетки нет; JCVI-syn3.0 и шкала Шарова в \S~\ref{sec:freeze} задают лишь \emph{грубую верхнюю} оценку. Пока C* не привязан к измеримой структуре, энергетические сравнения с пылинкой и хронология «заморозки» перед первой клеткой остаются качественными.

\subsection{«Заморозка» космической фазы на C*}
\label{sec:open_freeze}

Что именно останавливает дальнейшее усложнение выше C* в открытом космосе, не установлено. В \S~\ref{sec:freeze} принята рабочая гипотеза~--- барьер \emph{условий среды} (жидкофазная среда, ресурсный пул, длительный поток свободной энергии), а не глобального дефицита свободной энергии или «остывания» Вселенной; альтернативы (химический потолок репликации, неустойчивость протоклеток в вакууме, иные механизмы) не разведены.

\subsection{Носители пребиотической химии}
\label{sec:open_carriers}

Где в космической фазе накапливается сложность до порога C*~--- в молекулярных облаках, пыли (IDPs), на кометах или в метеоритах,~--- окончательно не установлено; вероятно, речь не об одном статичном носителе, а о \emph{цепочке} площадок. Чтобы ответить, необходимо проработать \emph{комплексный мультимодальный} процесс: накопление и перенос сложности между типами носителей с учётом \emph{времени жизни} каждого (угасание локальных источников энергии, разрушение зёрен, сублимация льда, абляция при доставке на планету) и \emph{смены носителей}~--- от молекулярных облаков через протопланетный диск к кометам, межпланетной пыли и, наконец, к планетной поверхности (\S~\ref{sec:idp_dust}, \S~\ref{sec:bimodal}). Пока эта цепочка не смоделирована количественно, разграничение между носителями накопления и носителем доставки остаётся качественным; рабочая гипотеза модели~--- \S~\ref{sec:idp_dust}.

\subsection{Конвергентность формы цивилизации}
\label{sec:open_civ_form}

Спектр «гуманоид функционального типа $\leftrightarrow$ \emph{Homo sapiens} вплотную» и границы, которые на него накладывает «трубка» $\sigma$, разобраны в \S~\ref{sec:time_form}. Кратко: из модели \emph{следует} \emph{функциональный} план гуманоидного носителя цивилизации (\S~\ref{sec:time_form}), но не утверждается совпадение с земной морфологией; открытый вопрос~--- \emph{где} на спектре окажется реальная форма: может быть близка к нам, а может сильно отличаться в морфологии и техносфере. Эмпирическая опора~--- лишь конвергентная эволюция (\S~\ref{sec:ltee}); фальсификация~--- \S~\ref{sec:pred_convergent_form}.

\subsection{Порог громкости и \texorpdfstring{$t_{\mathrm{loud}}$}{t\_loud}}
\label{sec:open_loud}

«Громкость» отождествлена с Типом~II по Кардашёву (\S~\ref{sec:loud_time}); порог заимствован из SETI-литературы~\cite{Wright2014,Hanson2021}. Открытый вопрос~--- насколько этот порог соответствует реальной детектируемости и применим ли тот же порог к Типу~I ($10^{16}$~Вт). Форвард-прогнозы: типовой подъём Тип~I\,$\to$\,II за $t_{\mathrm{loud}} \sim 10^3$~лет (\S~\ref{sec:loud_time}, \S~\ref{sec:pred_kardashev}) и условное число $\mathbb{E}[V_{\mathrm{I}}] \approx 5$ видимых соседей при гипотетической детектируемости Типа~I (подразд.~\ref{sec:pred_type1}); чувствительность $t_{\mathrm{loud}}$ и $t_{\mathrm{I}}$ к темпу роста энергопотребления~--- табл.~1.

\subsection{Галактическая конкуренция}
\label{sec:galactic_competition}

В модели $n \sim 10^4$--$10^5$ цивилизаций синхронизированной когорты в нашей галактике (\S~\ref{sec:kepler}). Открытый вопрос~--- как после Типа~II мы будем конкурировать за переход к Типу~III?

\bibliographystyle{unsrt}
\bibliography{refs}

\end{document}
